\documentclass[12pt,a4paper]{ctexart}
\usepackage{amsmath, amssymb, amsthm}
\usepackage{graphicx}
\usepackage{geometry}
\usepackage{tasks}
\usepackage{booktabs}
\usepackage{fancyhdr}
\usepackage{float}
\usepackage{mathrsfs}

% ==================================================
% 【核心修改】引入 hyperref 宏包以生成 PDF 书签
% ==================================================
\usepackage[
bookmarks=true,         % 开启书签
bookmarksnumbered=true, % 书签中包含章节编号 (如 "1. 引言" 而不是 "引言")
colorlinks=true,        % 开启超链接颜色 (设为 false 则文字为黑色，但可能有方框)
linkcolor=black,        % 目录/内部链接颜色 (设为 black 打印时更美观)
urlcolor=blue,          % 网址链接颜色
citecolor=green,        % 引用链接颜色
pdfstartview=FitH,      % 打开 PDF 时自动适合宽度
unicode=true            % 避免中文书签乱码 (对于 CTeX 通常自动识别，显式加上更稳妥)
]{hyperref}
% ==================================================

\newtheorem*{note}{注记}
% 定义定理环境，[section] 表示编号跟随章节（如 定理 2.1）
\newtheorem{theorem}{定理}[section]

% 定义引理环境，[theorem] 表示与定理共用编号计数（如 定理 2.2 之后是 引理 2.3）
\newtheorem{lemma}[theorem]{引理}
\newtheorem{definition}{定义}
\newtheorem{proposition}{命题}

\title{\textbf{重点展开}}
\author{Ao Yu}
\date{\today}
\ctexset{
	section = {
		format = \Large\bfseries\raggedright
	}
}

\begin{document}
	
	\maketitle
	
	% 可选：生成文档内部的目录
	% \tableofcontents 
	写在前面：
	\begin{align*}
		% 高斯公式 (Divergence Theorem)
		\text{高斯公式：} \quad & \int_{\Omega} \nabla \cdot \mathbf{F} \, \mathrm{d}V = \int_{\partial \Omega} \mathbf{F} \cdot \mathbf{n} \, \mathrm{d}S \\[10pt]
		% 格林第一公式 (Green's First Identity) -- 注意这里不能有空行
		\text{格林第一公式：} \quad & \int_{\Omega} \nabla u \cdot \nabla v \, \mathrm{d}V = -\int_{\Omega} v \Delta u \, \mathrm{d}V + \int_{\partial \Omega} v \frac{\partial u}{\partial n} \, \mathrm{d}S \\[10pt]
		% 格林第二公式 (Green's Second Identity) -- 这里也不能有空行
		\text{格林第二公式：} \quad & \int_{\Omega} (v \Delta u - u \Delta v) \, \mathrm{d}V = \int_{\partial \Omega} \left( v \frac{\partial u}{\partial n} - u \frac{\partial v}{\partial n} \right) \mathrm{d}S
	\end{align*}
	\newpage
	\section{引言}
	
	必考一题原
	
	\subsection{极小曲面问题}

	设 $\Omega \subset \mathbb{R}^2$ 为有界光滑区域，边界 $\partial\Omega$ 上给定边值
	\[
	u|_{\partial\Omega} = \varphi .
	\]
	在所有满足该边界条件的函数中，求一函数 $u(x,y)$，使其图像曲面
	\[
	S=\{(x,y,u(x,y)):(x,y)\in\Omega\}
	\]
	的表面积最小。
	
	\medskip
	
	\textbf{泛函形式}  
	定义允许函数集合
	\[
	M_\varphi=\{v\in C^1(\overline{\Omega}) : v|_{\partial\Omega}=\varphi\},
	\]
	表面积泛函为
	\[
	J(v)=\iint_{\Omega}\sqrt{1+v_x^2+v_y^2}\,dx\,dy.
	\]
	极小曲面问题等价于变分问题
	\[
	J(u)=\min_{v\in M_\varphi} J(v).
	\]
	
	\medskip
	
	\textbf{极值的必要条件（变分推导）}
	
	设 $u$ 为变分问题
	\[
	J(u)=\min_{v\in M_\varphi}J(v)
	\]
	的解，其中
	\[
	J(v)=\iint_{\Omega}\sqrt{1+v_x^2+v_y^2}\,dx\,dy.
	\]
	
	取任意变分函数
	\[
	v\in M_0:=\{v\in C^1(\overline{\Omega}): v|_{\partial\Omega}=0\},
	\]
	对任意 $\varepsilon\in\mathbb{R}$，有 $u+\varepsilon v\in M_\varphi$。定义
	\[
	j(\varepsilon)=J(u+\varepsilon v),
	\]
	则 $j(\varepsilon)$ 在 $\varepsilon=0$ 处取极小值，因此
	\[
	j'(0)=0.
	\]
	
	\medskip
	\textbf{一阶变分}
	
	直接计算得
	\[
	j'(\varepsilon)
	= \iint_{\Omega}
	\frac{(u_x+\varepsilon v_x)v_x+(u_y+\varepsilon v_y)v_y}
	{\sqrt{1+(u_x+\varepsilon v_x)^2+(u_y+\varepsilon v_y)^2}}
	\,dx\,dy.
	\]
	令 $\varepsilon=0$，得到
	\[
	\iint_{\Omega}
	\left(
	\frac{u_x}{\sqrt{1+u_x^2+u_y^2}}\,v_x
	+
	\frac{u_y}{\sqrt{1+u_x^2+u_y^2}}\,v_y
	\right)\,dx\,dy
	=0,
	\quad \forall v\in M_0.
	\]
	
	\medskip
	\textbf{Euler--Lagrange 方程}
	
	若 $u\in C^2(\Omega)$，由 Green 公式（$\iint_{\Omega} P \frac{\partial v}{\partial x} \, dx dy = \oint_{\partial \Omega} P v n_x \, ds - \iint_{\Omega} \frac{\partial P}{\partial x} v \, dx dy$）可得
	\[
	-\iint_{\Omega}
	\left[
	\frac{\partial}{\partial x}
	\left(
	\frac{u_x}{\sqrt{1+u_x^2+u_y^2}}
	\right)
	+
	\frac{\partial}{\partial y}
	\left(
	\frac{u_y}{\sqrt{1+u_x^2+u_y^2}}
	\right)
	\right]v\,dx\,dy
	=0.
	\]
	由于 $v$ 在 $\Omega$ 内任意，得到极小曲面方程
	\[
	\frac{\partial}{\partial x}
	\left(
	\frac{u_x}{\sqrt{1+u_x^2+u_y^2}}
	\right)
	+
	\frac{\partial}{\partial y}
	\left(
	\frac{u_y}{\sqrt{1+u_x^2+u_y^2}}
	\right)
	=0,
	\quad (x,y)\in\Omega.
	\]
	
	边界条件为
	\[
	u|_{\partial\Omega}=\varphi.
	\]
	
	\medskip
	\textbf{二阶变分与充分性}
	
	计算二阶变分得
	\[
	j''(\varepsilon)
	=
	\iint_{\Omega}
	\frac{
		v_x^2+v_y^2+
		\big[v_y(u_x+\varepsilon v_x)-v_x(u_y+\varepsilon v_y)\big]^2
	}
	{\big[1+(u_x+\varepsilon v_x)^2+(u_y+\varepsilon v_y)^2\big]^{3/2}}
	\,dx\,dy
	>0.
	\]
	因此，若边值问题
	\[
	\begin{cases}
		\displaystyle
		\frac{\partial}{\partial x}
		\left(
		\frac{u_x}{\sqrt{1+u_x^2+u_y^2}}
		\right)
		+
		\frac{\partial}{\partial y}
		\left(
		\frac{u_y}{\sqrt{1+u_x^2+u_y^2}}
		\right)=0,
		& \text{in }\Omega,\\[1ex]
		u=\varphi, & \text{on }\partial\Omega
	\end{cases}
	\]
	存在解 $u\in C^1(\overline{\Omega})\cap C^2(\Omega)$，
	则该解必为变分问题的极小解。
	
	\medskip
	\textbf{等价形式与线性化}
	
	极小曲面方程亦可写成
	\[
	(1+u_y^2)u_{xx}-2u_xu_yu_{xy}+(1+u_x^2)u_{yy}=0,
	\]
	这是一个非线性二阶椭圆方程。
	
	若假设
	\[
	|u_x|,\,|u_y|\ll1,
	\]
	忽略高阶小量，则方程线性化为
	\[
	u_{xx}+u_{yy}=0.
	\]
	因此，在小梯度假设下，极小曲面可近似看作 Laplace 方程
	第一类边值问题的解。
	
	\subsection{膜平衡问题}
	
	设薄膜占据 $Oxy$ 平面上的区域 $\Omega$，其边界由两部分组成：$\partial\Omega = \gamma \cup \Gamma$。
	
	\begin{itemize}
		\item 在 $\gamma$ 上：给定薄膜的位移（第一类/Dirichlet 边界）；
		\item 在 $\Gamma$ 上：给定作用在边界上的垂直外力密度 $p(x,y)$（第二类/Neumann 边界）。
	\end{itemize}
	此外，薄膜内部受到垂直方向的外力密度 $f(x,y)$ 作用。
	
	\medskip
	\textbf{泛函形式}
	
	根据最小势能原理，平衡状态对应于总势能的极小值。
	定义允许函数集合（满足强制边界条件）：
	\[
	M_\varphi = \{ v \in C^1(\overline{\Omega}) : v|_\gamma = \varphi \}.
	\]
	
	在小变形假设下（$|\nabla v| \ll 1$），总势能泛函 $J(v)$ 由变形能、外力势能和边界力势能组成：
	\[
	J(v) = \underbrace{\frac{T}{2} \iint_{\Omega} (v_x^2 + v_y^2) \,dx\,dy}_{\text{变形能}} 
	- \underbrace{\iint_{\Omega} f v \,dx\,dy}_{\text{外力做功}} 
	- \underbrace{\int_{\Gamma} p v \,ds}_{\text{边界力做功}},
	\]
	其中 $T$ 为薄膜张力常数。膜平衡问题即求解变分问题：
	\[
	J(u) = \min_{v \in M_\varphi} J(v).
	\]
	
	\medskip
	\textbf{极值的必要条件}
	
	设 $u \in M_\varphi$ 为使得总势能最小的解。取任意测试函数（满足零边值条件）：
	\[
	v \in M_0 := \{ v \in C^1(\overline{\Omega}) : v|_\gamma = 0 \}.
	\]
	对任意实数 $\varepsilon$，函数 $j(\varepsilon) = J(u + \varepsilon v)$ 在 $\varepsilon = 0$ 处取得极小值，故 $j'(0) = 0$。
	
	\medskip
	\textbf{一阶变分（弱形式）}
	
	计算泛函的一阶变分：
	\[
	\delta J = \frac{d}{d\varepsilon} J(u+\varepsilon v) \bigg|_{\varepsilon=0}
	= T \iint_{\Omega} (\nabla u \cdot \nabla v) \,dx\,dy - \iint_{\Omega} f v \,dx\,dy - \int_{\Gamma} p v \,ds = 0.
	\]
	由此得到问题的\textbf{弱形式}：寻找 $u \in M_\varphi$，使得
	\[
	T \iint_{\Omega} \nabla u \cdot \nabla v \,dx\,dy = \iint_{\Omega} f v \,dx\,dy + \int_{\Gamma} p v \,ds, \quad \forall v \in M_0.
	\]
	
	\medskip
	\textbf{Euler--Lagrange 方程（强形式）}
	
	若解 $u$ 足够光滑（例如 $u \in C^2(\Omega)$），应用 Green 第一公式（分部积分）：
	\[
	\iint_{\Omega} \nabla u \cdot \nabla v \,dx\,dy = - \iint_{\Omega} (\Delta u) v \,dx\,dy + \int_{\partial\Omega} \frac{\partial u}{\partial n} v \,ds.
	\]
	将此代入弱形式方程中。注意边界积分分为 $\gamma$ 和 $\Gamma$ 两部分，且在 $\gamma$ 上 $v=0$，因此沿 $\gamma$ 的积分项消失：
	\[
	-\iint_{\Omega} (T \Delta u + f) v \,dx\,dy + \int_{\Gamma} \left( T \frac{\partial u}{\partial n} - p \right) v \,ds = 0.
	\]
	
	根据变分法基本引理，由于 $v$ 在区域 $\Omega$ 内部任意，且在边界 $\Gamma$ 上也任意，上述等式成立的充要条件是各项被积函数分别为零。由此导出方程及边界条件：
	
	\[
	\begin{cases}
		-T \Delta u = f, & (x,y) \in \Omega \quad (\text{Poisson 方程}), \\
		u = \varphi, & (x,y) \in \gamma \quad (\text{强制边界条件}), \\
		\displaystyle T \frac{\partial u}{\partial n} = p, & (x,y) \in \Gamma \quad (\text{自然边界条件}).
	\end{cases}
	\]
	
	\medskip
	\textbf{注记}：
	与第一类边界条件（$u|_\gamma=\varphi$）不同，第二类边界条件（$T \frac{\partial u}{\partial n}|_\Gamma = p$）并非预先强加在函数空间 $M_\varphi$ 中，而是由变分过程自然导出的，因此称为\textbf{自然边界条件}。
	\vspace{3cm}
	
	\textbf{例：}
	设
	\[
	J(v) = \frac{1}{2}\int_{\Omega} (|\nabla v|^2 + v^2) \,\mathrm{d}x + \frac{1}{2}\int_{\partial\Omega} \alpha(x)v^2 \,\mathrm{d}s - \int_{\Omega} fv \,\mathrm{d}x - \int_{\partial\Omega} gv \,\mathrm{d}s,
	\]
	其中 $\alpha(x) \geqslant 0$。考虑以下三个问题：
	
	\noindent\textbf{问题 I (变分问题)：} 求 $u \in M = C^1(\overline{\Omega})$，使得
	\[
	J(u) = \min_{v \in M} J(v).
	\]
	
	\noindent\textbf{问题 II：} 求 $u \in M = C^1(\overline{\Omega})$，使得它对于任意 $v \in M$，都满足
	\[
	\int_{\Omega} (\nabla u \cdot \nabla v + u \cdot v - fv) \,\mathrm{d}x + \int_{\partial\Omega} (\alpha(x)uv - gv) \,\mathrm{d}s = 0.
	\]
	
	\noindent\textbf{问题 III (第三边值问题)：} 求 $u \in C^2(\Omega) \cap C^1(\overline{\Omega})$，满足以下边值问题
	\[
	\begin{cases}
		-\Delta u + u = f, & x \in \Omega, \\
		\dfrac{\partial u}{\partial n} + \alpha(x)u = g, & x \in \partial\Omega.
	\end{cases}
	\]
	
	\begin{enumerate}
		\item[(1)] 证明问题 I 与问题 II 等价。
		\item[(2)] 当 $u \in C^2(\Omega) \cap C^1(\overline{\Omega})$ 时，证明问题 I、II、III 等价。
	\end{enumerate}
	
	\noindent\textbf{证明}
	
	\noindent\fbox{
		\parbox{0.95\textwidth}{
			\textbf{补充知识 (Green 第一恒等式)：}
			\[
			\int_{\Omega} \nabla u \cdot \nabla v \,\mathrm{d}x = -\int_{\Omega} (\Delta u)v \,\mathrm{d}x + \int_{\partial\Omega} \frac{\partial u}{\partial n}v \,\mathrm{d}s.
			\]
		}
	}
	
	\begin{proof}
		\textbf{(1) 证明问题 I 与问题 II 等价}
		
		\noindent ($\Longrightarrow$) \textbf{必要性：}
		取 $u = \arg\min_{v \in M} J(v)$ 且 $u \in M$。对任意 $v \in M$，定义 $j(\varepsilon) = J(u + \varepsilon v)$。则我们应保证 $j'(0) = 0$ 且 $j''(0) \geqslant 0$。
		
		展开 $j(\varepsilon)$：
		\[
		\begin{aligned}
			j(\varepsilon) &= \frac{1}{2}\int_{\Omega} [|\nabla(u+\varepsilon v)|^2 + (u+\varepsilon v)^2] \,\mathrm{d}x + \frac{1}{2}\int_{\partial\Omega} \alpha(u+\varepsilon v)^2 \,\mathrm{d}s \\
			&\quad - \int_{\Omega} f(u+\varepsilon v) \,\mathrm{d}x - \int_{\partial\Omega} g(u+\varepsilon v) \,\mathrm{d}s.
		\end{aligned}
		\]
		对 $\varepsilon$ 求导并令 $\varepsilon = 0$：
		\[
		j'(0) = \int_{\Omega} (\nabla u \cdot \nabla v + uv) \,\mathrm{d}x + \int_{\partial\Omega} \alpha uv \,\mathrm{d}s - \int_{\Omega} fv \,\mathrm{d}x - \int_{\partial\Omega} gv \,\mathrm{d}s.
		\]
		令 $j'(0)=0$，整理即得：
		\[
		\int_{\Omega} (\nabla u \cdot \nabla v + uv - fv) \,\mathrm{d}x + \int_{\partial\Omega} (\alpha uv - gv) \,\mathrm{d}s = 0.
		\]
		此即问题 II 的方程。
		
		\vspace{1em}
		\noindent ($\Longleftarrow$) \textbf{充分性：}
		为了简化书写，定义双线性形式 $a(u,v)$ 和线性泛函 $l(v)$：
		\[
		a(u,v) = \int_{\Omega} (\nabla u \cdot \nabla v + uv) \,\mathrm{d}x + \int_{\partial\Omega} \alpha uv \,\mathrm{d}s, \quad l(v) = \int_{\Omega} fv \,\mathrm{d}x + \int_{\partial\Omega} gv \,\mathrm{d}s.
		\]
		则泛函可重写为 $J(v) = \frac{1}{2}a(v,v) - l(v)$。
		
		已知 $u$ 满足问题 II，即对任意 $v \in M$，有 $a(u,v) = l(v)$ 成立。
		我们取任意 $w \in M$，令 $w = u - v$ (即 $v = u - w$，注意此处 $w$ 为偏差量)，考察 $J(v) - J(u)$：
		\[
		\begin{aligned}
			J(v) - J(u) &= \left[ \frac{1}{2}a(u-w, u-w) - l(u-w) \right] - \left[ \frac{1}{2}a(u,u) - l(u) \right] \\
			&= \frac{1}{2}[a(u,u) - 2a(u,w) + a(w,w)] - l(u) + l(w) - \frac{1}{2}a(u,u) + l(u) \\
			&= \frac{1}{2}a(w,w) - [a(u,w) - l(w)].
		\end{aligned}
		\]
		由于 $u$ 满足 $a(u,w) = l(w)$，故中括号项为 0。于是：
		\[
		J(v) - J(u) = \frac{1}{2}a(w,w) = \frac{1}{2}\int_{\Omega} (|\nabla w|^2 + w^2) \,\mathrm{d}x + \frac{1}{2}\int_{\partial\Omega} \alpha w^2 \,\mathrm{d}s \geqslant 0.
		\]
		当且仅当 $w=0$ 即 $v=u$ 时取等号。故 $u$ 是唯一极小值点。
		
		\vspace{1em}
		\textbf{(2) 当 $u \in C^2(\Omega) \cap C^1(\overline{\Omega})$ 时，证明问题 I、II、III 等价}
		
		由 (1) 知 I $\Leftrightarrow$ II，故只需证 II $\Leftrightarrow$ III。
		
		\noindent ($\Longrightarrow$) \textbf{II 推导 III：}
		结合 Green 第一恒等式，我们将问题 II 的等式变形：
		\[
		\int_{\Omega} (-\Delta u + u - f)v \,\mathrm{d}x + \int_{\partial\Omega} \left( \frac{\partial u}{\partial n} + \alpha u - g \right)v \,\mathrm{d}s = 0.
		\]
		结合 $v$ 的任意性（变分法基本引理）：
		\begin{itemize}
			\item 取 $v \in C_0^{\infty}(\Omega)$，可得 $-\Delta u + u = f, \quad x \in \Omega$。
			\item 代回原式，余下边界项需为 0，可得 $\dfrac{\partial u}{\partial n} + \alpha u = g, \quad x \in \partial\Omega$。
		\end{itemize}
		
		\noindent ($\Longleftarrow$) \textbf{III 推导 II：}
		结合 $-\Delta u + u = f, \, x \in \Omega$，对两边同时乘以 $v$ 并积分：
		\[
		\int_{\Omega} (-\Delta u + u - f)v \,\mathrm{d}x = 0.
		\]
		同理，利用边界条件 $\frac{\partial u}{\partial n} + \alpha u - g = 0$，在边界积分：
		\[
		\int_{\partial\Omega} \left( \frac{\partial u}{\partial n} + \alpha u - g \right)v \,\mathrm{d}s = 0.
		\]
		将两者相加，并对第一项使用 Green 第一恒等式展开 $-\Delta u$ 项，即可还原得到：
		\[
		\int_{\Omega} (\nabla u \cdot \nabla v + u \cdot v - fv) \,\mathrm{d}x + \int_{\partial\Omega} (\alpha uv - gv) \,\mathrm{d}s = 0.
		\]
	\end{proof}
	
	\section{波动方程}
	
	\subsection{特征线法求解一阶线性微分方程}
	
	考一题计算
	
	我们称下列常微分方程初值问题
	\begin{equation}
		\begin{cases}
			\dfrac{dx}{dt} = a \\[6pt]
			x(0) = c
		\end{cases}
	\end{equation}
	的解 $x(t,c) = at + c$ 为方程（Cauchy 问题）
	\begin{equation}
		\frac{\partial \rho}{\partial t} + a \frac{\partial \rho}{\partial x} = 0
	\end{equation}
	的特征线，其中 $c$ 为常数。
	
	我们直接来考虑最复杂的情形：
	\begin{equation}
		\begin{cases}
			\dfrac{\partial u}{\partial t} + f(x,t)\dfrac{\partial u}{\partial x} = g(u; x, t), & x \in \mathbb{R}^1, t > 0 \\[8pt]
			u(x,0) = h(x).
		\end{cases}
	\end{equation}
	
	\textbf{解：} 我们构造特征线
	\begin{equation}
		\begin{cases}
			\dfrac{dx}{dt} = f(x,t) \\[6pt]
			x(0) = c
		\end{cases}
	\end{equation}
	得到方程的解为 $x(t,c)$，其中 $c$ 可以由 $x,t$ 表示：$c = c_0(x,t)$。
	
	取 $U(t) = u(x(t,c), t)$，则我们有
	\begin{equation}
		\begin{cases}
			\dfrac{dU}{dt} = \dfrac{\partial u}{\partial x} \dfrac{dx}{dt} + \dfrac{\partial u}{\partial t} = \dfrac{\partial u}{\partial t} + f(x,t)\dfrac{\partial u}{\partial x} = g(U; x(t,c), t) \\[10pt]
			U(0) = u(x(0,c), 0) = u(c,0) = h(c).
		\end{cases}
	\end{equation}
	此时我们将问题转化为求解上述 ODE 问题，得到方程的解为 $U(t,c) = U(t, c_0(x,t))$。
	
	另外，我们也可以考虑如下问题，它们是原问题的简化形式：
	\begin{equation}
		\frac{\partial u}{\partial t} + a \frac{\partial u}{\partial x} = 0 \quad \bigg/ \quad \frac{\partial u}{\partial t} + \frac{\partial}{\partial x}(q(u(x))) = 0
	\end{equation}
	
	\subsubsection{习题练习}
	\textbf{例：} 求下列 Cauchy 问题的解
	\begin{equation}
		\begin{cases}
			\dfrac{\partial u}{\partial t} + (x+t)\dfrac{\partial u}{\partial x} + u = x, & x \in \mathbb{R}^1, t > 0 \\[8pt]
			u(x,0) = x.
		\end{cases} \label{eq:ex_pde}
	\end{equation}
	
	\textbf{解：}
	
	\textbf{第一步：} 求特征线、特征方程
	\begin{equation}
		\begin{cases}
			\dfrac{dx}{dt} = x + t \\[6pt]
			x(0) = c
		\end{cases}
	\end{equation}
	这是一个一阶线性非齐次常微分方程，其解为：
	\begin{equation}
		x(t) = e^t(1+c) - (1+t). \label{eq:ex_char_sol}
	\end{equation}
	
	\textbf{第二步：} 令 $U(t) = u(x(t), t)$，则由 \eqref{eq:ex_pde}, \eqref{eq:ex_char_sol} 可得 ODE 初值问题：
	\begin{equation}
		\begin{cases}
			\dfrac{dU(t)}{dt} + U(t) = e^t(1+c) - (1+t), \\[8pt]
			U(0) = u(x(0),0) = u(c,0) = c.
		\end{cases}
	\end{equation}
	求解该一阶线性微分方程，其解为：
	\begin{equation}
		U(t) = \frac{1}{2}(c+1)e^t + \frac{1}{2}(c-1)e^{-t} - t. \label{eq:ex_U_sol}
	\end{equation}
	
	\textbf{第三步：} 从 \eqref{eq:ex_char_sol} 中解出 $c$：
	\begin{equation}
		c = (x+t+1)e^{-t} - 1,
	\end{equation}
	然后代入 \eqref{eq:ex_U_sol} 得所求的解为：
	\begin{equation}
		u(x,t) = \frac{1}{2}e^{-2t}(x+t+1) - e^{-t} + \frac{1}{2}(x-t+1).
	\end{equation}

	
	\textbf{例：} 求解下列初值问题
	\begin{equation}
		\begin{cases}
			u_t + 2u_x = 0, & t > 0, -\infty < x < \infty, \\[6pt]
			u|_{t=0} = x^2, & -\infty < x < \infty.
		\end{cases}
	\end{equation}
	
	\textbf{解：}
	特征方程为 $\dfrac{dx}{dt} = 2$，初始条件 $x(0)=c$。
	解得特征线为 $x = 2t + c$，即 $c = x - 2t$。
	
	沿特征线，方程化为 $\dfrac{du}{dt} = 0$，故 $u$ 沿特征线为常数。
	由初始条件得：
	\[ U(t) = u(x(0), 0) = u(c, 0) = c^2. \]
	将 $c = x - 2t$ 代入，得原方程的解为：
	\begin{equation}
		u(x,t) = (x - 2t)^2.
	\end{equation}
	
	\vspace{1em} % 增加一点垂直间距
	
	\textbf{例：} 求解下列初值问题
	\begin{equation}
		\begin{cases}
			u_t + (1+x^2)u_x - u = 0, & t > 0, -\infty < x < \infty, \\[6pt]
			u|_{t=0} = \arctan x, & -\infty < x < \infty.
		\end{cases}
	\end{equation}
	
	\textbf{解：}
	方程可改写为 $u_t + (1+x^2)u_x = u$。
	特征方程为：
	\begin{equation}
		\begin{cases}
			\dfrac{dx}{dt} = 1 + x^2 \\[6pt]
			x(0) = c
		\end{cases}
	\end{equation}
	分离变量积分 $\displaystyle\int \frac{dx}{1+x^2} = \int dt$，得 $\arctan x = t + K$。
	代入初始条件 $t=0, x=c$，得 $K = \arctan c$。
	故特征关系式为：
	\[ \arctan x = t + \arctan c \implies \arctan c = \arctan x - t. \]
	
	沿特征线，原方程化为 $\dfrac{du}{dt} = u$，初始条件 $U(0) = u(c,0) = \arctan c$。
	解该 ODE 得：
	\[ U(t) = (\arctan c) e^t. \]
	将 $\arctan c = \arctan x - t$ 代入，得原方程的解为：
	\begin{equation}
		u(x,t) = (\arctan x - t)e^t.
	\end{equation}
	
	\subsection{球面平均法与 Kirchhoff 公式的推导}

	必考推导过程

	我们在三维空间 $\mathbb{R}^3 \times (0, \infty)$ 中考虑如下的柯西问题（初值问题）：
	\begin{equation}
		\begin{cases}
			\dfrac{\partial^2 u}{\partial t^2} - a^2 \Delta u = f(x,t), & x \in \mathbb{R}^3, t > 0 \\
			u|_{t=0} = \varphi(x), & x \in \mathbb{R}^3 \\
			u_t|_{t=0} = \psi(x), & x \in \mathbb{R}^3
		\end{cases}
		\label{eq:wave_3d}
	\end{equation}
	利用线性叠加原理，我们只需分别考虑 $f=\varphi=0$（仅 $\psi \neq 0$ 的情形），然后通过关于 $t$ 的求导和 Duhamel 原理即可得到一般解。
	
	\subsubsection*{1. 球面平均函数的定义及其性质}
	
	对任意连续函数 $h(x) \in C(\mathbb{R}^3)$，定义其在以 $x$ 为心、$r$ 为半径的球面上的平均值为：
	\begin{equation}
		I(x,r;h) = \dfrac{1}{4\pi r^2} \iint_{|y-x|=r} h(y) \, \mathrm{d}S_y = \dfrac{1}{4\pi} \iint_{|\eta|=1} h(x+r\eta) \, \mathrm{d}S_\eta
		\label{eq:spherical_mean_def}
	\end{equation}
	我们首先建立 $I(x,r;h)$ 满足的微分方程。这是一个关键步骤，它联系了拉普拉斯算子与径向导数。
	
	考虑球体上的积分，利用高斯公式（散度定理）：
	\begin{align*}
		\iiint_{|z| \le r} \Delta_z h(x+z) \, \mathrm{d}z &= \iint_{|z|=r} \dfrac{\partial h}{\partial n} \, \mathrm{d}S_z \\
		&= \iint_{|z|=r} \sum_{i=1}^3 \dfrac{\partial h}{\partial z_i} \cdot \dfrac{z_i}{r} \, \mathrm{d}S_z \\
		&= r^2 \iint_{|\eta|=1} \sum_{i=1}^3 \dfrac{\partial h}{\partial x_i}(x+r\eta) \cdot \eta_i \, \mathrm{d}S_\eta \quad (\text{令 } z=r\eta) \\
		&= r^2 \dfrac{\partial}{\partial r} \iint_{|\eta|=1} h(x+r\eta) \, \mathrm{d}S_\eta \\
		&= r^2 \dfrac{\partial}{\partial r} \left( 4\pi I(x,r;h) \right) \\
		&= 4\pi r^2 \dfrac{\partial I}{\partial r}
	\end{align*}
	另一方面，利用拉普拉斯算子的平移不变性 $\Delta_z h(x+z) = \Delta_x h(x+z)$，并在球坐标下积分：
	\begin{align*}
		\iiint_{|z| \le r} \Delta_z h(x+z) \, \mathrm{d}z &= \iiint_{|z| \le r} \Delta_x h(x+z) \, \mathrm{d}z \\
		&= \Delta_x \iiint_{|z| \le r} h(x+z) \, \mathrm{d}z \\
		&= \Delta_x \int_0^r \rho^2 \left( \iint_{|\eta|=1} h(x+\rho\eta) \, \mathrm{d}S_\eta \right) \mathrm{d}\rho \\
		&= \Delta_x \int_0^r 4\pi \rho^2 I(x,\rho;h) \, \mathrm{d}\rho
	\end{align*}
	对上述两式关于 $r$ 求导（利用变上限积分求导法则）：
	\[
	\dfrac{\partial}{\partial r} \left( 4\pi r^2 \dfrac{\partial I}{\partial r} \right) = \dfrac{\partial}{\partial r} \left( \Delta_x \int_0^r 4\pi \rho^2 I(x,\rho;h) \, \mathrm{d}\rho \right)
	\]
	\[
	4\pi \dfrac{\partial}{\partial r} \left( r^2 \dfrac{\partial I}{\partial r} \right) = 4\pi r^2 \Delta_x I(x,r;h)
	\]
	整理得到著名的Darboux 方程：
	\begin{equation}
		\Delta_x I = \dfrac{1}{r^2} \dfrac{\partial}{\partial r} \left( r^2 \dfrac{\partial I}{\partial r} \right) = \dfrac{\partial^2 I}{\partial r^2} + \dfrac{2}{r} \dfrac{\partial I}{\partial r}
		\label{eq:darboux}
	\end{equation}
	
	\subsubsection*{2. 辅助函数 $M(x,r,t)$ 的引入与降维}
	
	现在回到波动方程 $u_{tt} = a^2 \Delta u$。对该方程两边取球面平均 $I(x,r;\cdot)$。
	由于 $t$ 与积分变量无关，且 $\Delta_x$ 与积分是线性的，我们有：
	\[
	\dfrac{\partial^2}{\partial t^2} I(x,r;u) = a^2 I(x,r;\Delta u) = a^2 \Delta_x I(x,r;u)
	\]
	利用 (\ref{eq:darboux}) 式，上述方程变为：
	\[
	\dfrac{\partial^2 I}{\partial t^2} = a^2 \left( \dfrac{\partial^2 I}{\partial r^2} + \dfrac{2}{r} \dfrac{\partial I}{\partial r} \right)
	\]
	这是一个带有奇异系数 $\frac{2}{r}$ 的方程。为了消去一阶导数项，我们引入变换：
	\begin{equation}
		M(x,r,t) = r I(x,r;u)
	\end{equation}
	计算 $M$ 对 $r$ 的导数：
	\[
	\dfrac{\partial M}{\partial r} = I + r \dfrac{\partial I}{\partial r}, \quad \dfrac{\partial^2 M}{\partial r^2} = 2\dfrac{\partial I}{\partial r} + r \dfrac{\partial^2 I}{\partial r^2} = r \left( \dfrac{\partial^2 I}{\partial r^2} + \dfrac{2}{r}\dfrac{\partial I}{\partial r} \right)
	\]
	因此，方程化为标准的一维波动方程：
	\begin{equation}
		\dfrac{\partial^2 M}{\partial t^2} = r \dfrac{\partial^2 I}{\partial t^2} = r \cdot a^2 \left( \dfrac{\partial^2 I}{\partial r^2} + \dfrac{2}{r} \dfrac{\partial I}{\partial r} \right) = a^2 \dfrac{\partial^2 M}{\partial r^2}
	\end{equation}
	这就把三维问题转化为了半无界区间 $r \in [0, \infty)$ 上的一维问题。
	
	\subsubsection*{3. 求解定解问题}
	
	首先考虑 $f = \varphi = 0, \psi \neq 0$ 的情形。此时 $M(x,r,t)$ 满足：
	\begin{equation}
		\begin{cases}
			M_{tt} = a^2 M_{rr}, & r > 0, t > 0 \\
			M|_{t=0} = r I(x,r;0) = 0, & r \ge 0 \\
			M_t|_{t=0} = r I(x,r;\psi), & r \ge 0 \\
			M|_{r=0} = 0 \cdot I(x,0;u) = 0, & t \ge 0 \quad (\text{边界条件})
		\end{cases}
	\end{equation}
	这是一个半无界弦振动问题（固定端点）。其解可以通过奇延拓转化为全直线上的 d'Alembert 公式。当 $0 \le r \le at$ 时，解为：
	\begin{equation}
		M(x,r,t) = \dfrac{1}{2a} \int_{at-r}^{at+r} \xi I(x,\xi;\psi) \, \mathrm{d}\xi
		\label{eq:M_solution}
	\end{equation}
	我们要找的是原函数 $u(x,t)$ 在 $x$ 处的值，这对应于 $r \to 0$ 的极限：
	\[
	u(x,t) = \lim_{r \to 0} I(x,r;u) = \lim_{r \to 0} \dfrac{M(x,r,t)}{r}
	\]
	由于 $M(x,0,t)=0$，利用导数定义（或 L'H\^{o}pital 法则）：
	\begin{align*}
		u(x,t) &= \left. \dfrac{\partial M(x,r,t)}{\partial r} \right|_{r=0} \\
		&= \dfrac{1}{2a} \dfrac{\partial}{\partial r} \left( \int_{at-r}^{at+r} \xi I(x,\xi;\psi) \, \mathrm{d}\xi \right) \bigg|_{r=0} \\
		&= \dfrac{1}{2a} \left[ (at+r)I(x,at+r;\psi) \cdot 1 - (at-r)I(x,at-r;\psi) \cdot (-1) \right] \bigg|_{r=0} \\
		&= \dfrac{1}{2a} \left[ 2at I(x,at;\psi) \right] \\
		&= t I(x,at;\psi)
	\end{align*}
	展开球面平均的定义，得到 $\varphi=f=0$ 时的解 $u_\psi$：
	\begin{equation}
		u_\psi(x,t) = t \cdot \dfrac{1}{4\pi} \iint_{|\eta|=1} \psi(x+at\eta) \, \mathrm{d}S_\eta = \dfrac{1}{4\pi a^2 t} \iint_{S_{at}(x)} \psi(y) \, \mathrm{d}S_y
	\end{equation}
	其中 $S_{at}(x)$ 表示以 $x$ 为心，$at$ 为半径的球面。
	
	\subsubsection*{4. Kirchhoff 公式的综合形式}
	
	(1) 对于 $\varphi \neq 0$ 的项：
	设 $v$ 是对应初值 $\psi=\varphi, \varphi=0$ 的解，则 $u_\varphi = \dfrac{\partial v}{\partial t}$ 是对应初值 $\varphi \neq 0, \psi=0$ 的解（Stokes 法则）。因此：
	\[
	u_\varphi(x,t) = \dfrac{\partial}{\partial t} \left( t I(x,at;\varphi) \right)
	\]
	
	(2) 对于非齐次项 $f(x,t)$：
	利用 Duhamel 原理，非齐次方程的解是齐次方程解对时间的积分：
	\[
	u_f(x,t) = \int_0^t \left[ (t-\tau) I(x, a(t-\tau); f(\cdot, \tau)) \right] \, \mathrm{d}\tau
	\]
	
	(3) 最终公式：
	将上述三部分相加，得到三维波动方程的 Kirchhoff 公式：
	\begin{align}
		u(x,t) &= \dfrac{\partial}{\partial t} \left( \dfrac{1}{4\pi a^2 t} \iint_{S_{at}(x)} \varphi(y) \, \mathrm{d}S_y \right) \nonumber \\
		&\quad + \dfrac{1}{4\pi a^2 t} \iint_{S_{at}(x)} \psi(y) \, \mathrm{d}S_y \nonumber \\
		&\quad + \int_0^t \dfrac{1}{4\pi a^2 (t-\tau)} \iint_{S_{a(t-\tau)}(x)} f(y,\tau) \, \mathrm{d}S_y \, \mathrm{d}\tau
		\label{eq:kirchhoff_final}
	\end{align}
	该公式清楚地表明了三维波动的惠更斯原理（Huygens' Principle）：$t$ 时刻在 $x$ 点的波扰动仅依赖于 $t=0$ 时刻以 $x$ 为心、半径为 $at$ 的球面上的数据（以及该球面上过去源项的积分），而与球内部的数据无关（无后效现象）。
	
	\subsection{特征值问题与分离变量}
	
	% 内容来自 image_6c3784.png
	\subsubsection{理论基础：Sturm-Liouville 特征值问题}
	\textbf{定理 2.6.1}：在 $[0, l]$ 上特征值问题
	\begin{equation}
		\begin{cases}
			X'' + \lambda X = 0, & 0 < x < l \\
			-\alpha_1 X'(0) + \beta_1 X(0) = 0 \\
			\alpha_2 X'(l) + \beta_2 X(l) = 0
		\end{cases}
	\end{equation}
	(其中 $\alpha_i \ge 0, \beta_i \ge 0, \alpha_i + \beta_i \neq 0$) 具有如下性质：
	
	\begin{enumerate}
		\item[(a)] 所有特征值都是非负实数，特别当 $\beta_1 + \beta_2 > 0$ 时，所有特征值都是正数。
		\item[(b)] 所有特征值组成一个单调递增以无穷远点为凝聚点的序列：
		\[
		0 \le \lambda_1 < \dots < \lambda_n < \dots, \quad \lim_{n\to\infty} \lambda_n = \infty
		\]
		\item[(c)] 不同特征值对应的特征函数必正交；即若 $X_\lambda(x)$ 和 $X_\mu(x)$ 分别对应于特征值 $\lambda, \mu (\lambda \neq \mu)$ 的特征函数，则
		\[
		\int_0^l X_\lambda(x) X_\mu(x) dx = 0
		\]
		\item[(d)] 任意函数 $f(x) \in L^2[0, l]$ 可以按特征函数系 $\{X_n(x)\}$ 展开，即
		\[
		f(x) = \sum_{n=1}^{\infty} C_n X_n(x)
		\]
		其中
		\[
		C_n = \frac{\int_0^l f(x)X_n(x)dx}{\int_0^l X_n^2(x)dx}
		\]
	\end{enumerate}
	
	% 内容来自 image_6c3ab0.png
	\subsubsection{典型例题：Neumann 边界条件}
	\textbf{例 2.6.1}：求以下特征值问题的特征函数
	\begin{equation}
		\begin{cases}
			X''(x) + \lambda X(x) = 0 \\
			X'(0) = X'(l) = 0
		\end{cases}
	\end{equation}
	
	\textbf{解答}：结合特征方程 $r^2 + \lambda = 0$，我们有
	\begin{align*}
		X(x) &= C_1 \sin \sqrt{\lambda}x + C_2 \cos \sqrt{\lambda}x \\
		X'(x) &= C_1 \sqrt{\lambda} \cos \sqrt{\lambda}x - C_2 \sqrt{\lambda} \sin \sqrt{\lambda}x
	\end{align*}
	代入边界条件：
	\begin{equation}
		\begin{cases}
			X'(0) = C_1 \sqrt{\lambda} = 0 \\
			X'(l) = C_1 \sqrt{\lambda} \cos \sqrt{\lambda}l - C_2 \sqrt{\lambda} \sin \sqrt{\lambda}l = 0
		\end{cases}
		\implies
		\begin{cases}
			C_1 = 0 \\
			\sqrt{\lambda} = \frac{n\pi}{l}, \quad n=1, 2, \dots
		\end{cases}
	\end{equation}
	则特征函数为 $X_n(x) = C \sin \frac{n\pi}{l}x, \quad n=1, 2, \dots$ 
	% 注：原文图片此处结果写为 sin，但根据 Neumann 边界条件 X'(0)=0 推导通常应为 cos，此处忠实还原图片内容。
	
	% 内容来自 image_6c3b24.png 和 image_6c3b47.png
	\subsubsection{波动方程混合边值问题求解}
	\textbf{题目}：求解如下定解问题
	\begin{equation}
		\begin{cases}
			Lu = 0, & (x, t) \in Q \\
			u|_{x=0} = u_x|_{x=l} = 0, & t \ge 0 \\
			u|_{t=0} = x(x-2l), \quad u_t|_{t=0} = 0, & 0 \le x \le l
		\end{cases}
	\end{equation}
	
	\textbf{解答}：
	由分离变量法，设 $u(x,t) = X(x)T(t)$。
	关于空间变量 $X(x)$ 满足方程 $X''(x) + \lambda X(x) = 0$，其通解为：
	\begin{align*}
		X(x) &= C_1 \sin \sqrt{\lambda}x + C_2 \cos \sqrt{\lambda}x \\
		X'(x) &= C_1 \sqrt{\lambda} \cos \sqrt{\lambda}x - C_2 \sqrt{\lambda} \sin \sqrt{\lambda}x
	\end{align*}
	结合边界条件：
	\begin{equation}
		\begin{cases}
			X(0) = C_2 = 0 \\
			X'(l) = C_1 \sqrt{\lambda} \cos \sqrt{\lambda}l - 0 = 0
		\end{cases}
		\implies
		\begin{cases}
			\sqrt{\lambda} = \frac{(2n+1)\pi}{2l}, \quad n=0, 1, \dots \\
			C_2 = 0
		\end{cases}
	\end{equation}
	则特征值 $\lambda_n = \left(\frac{(2n+1)\pi}{2l}\right)^2$，特征函数为 $X_n(x) = \sin \frac{(2n+1)\pi}{2l}x, \quad n=0, 1, \dots$。
	
	同理，时间部分 $T(t)$ 满足 $T'' + a^2\lambda T = 0$，考虑到初始速度 $u_t|_{t=0} = 0$，即 $T'(0)=0$，故只保留余弦项：
	\[
	T_n(t) = B_n \cos a\sqrt{\lambda_n}t = B_n \cos \frac{(2n+1)\pi a}{2l}t
	\]
	于是解的形式为级数叠加：
	\[
	u(x, t) = \sum_{n=0}^{\infty} B_n \cos \frac{(2n+1)\pi a}{2l}t \sin \frac{(2n+1)\pi}{2l}x
	\]
	利用初始位移条件确定系数：
	由 $u|_{t=0} = x(x - 2l)$，我们有：
	\[
	\sum_{n=0}^{\infty} B_n \sin \frac{(2n+1)\pi}{2l}x = x^2 - 2lx
	\]
	其中系数 $B_n$ 由广义傅里叶系数公式求得（利用正交性）：
	\[
	B_n = \frac{\int_0^l (x^2-2lx) \sin \frac{(2n+1)\pi}{2l}x \, dx}{\int_0^l \sin^2 \frac{(2n+1)\pi}{2l}x \, dx}
	\]
	分母积分为 $\frac{l}{2}$。分子部分记为 $I$，经分部积分计算得：
	\begin{align*}
		I &= \int_0^l (x^2-2lx) \sin \beta_n x \, dx \quad (\text{令 } \beta_n = \frac{(2n+1)\pi}{2l}) \\
		&= \left[ -\frac{x^2-2lx}{\beta_n} \cos \beta_n x \right]_0^l + \frac{1}{\beta_n} \int_0^l (2x-2l) \cos \beta_n x \, dx \\
		&= 0 + \frac{1}{\beta_n} \left( \left[ \frac{2x-2l}{\beta_n} \sin \beta_n x \right]_0^l - \int_0^l \frac{2}{\beta_n} \sin \beta_n x \, dx \right) \\
		&= -\frac{2}{\beta_n^2} \left[ -\frac{1}{\beta_n} \cos \beta_n x \right]_0^l = \frac{2}{\beta_n^3} (\cos \beta_n l - 1) = -\frac{2}{\beta_n^3}
	\end{align*}
	代回 $B_n$ 的表达式：
	\[
	B_n = \frac{-2/\beta_n^3}{l/2} = -\frac{4}{l \beta_n^3} = -\frac{4}{l} \left( \frac{2l}{(2n+1)\pi} \right)^3 = -\frac{32l^2}{(2n+1)^3\pi^3}
	\]
	因此，最终解为：
	\begin{equation}
		u(x, t) = -\frac{32l^2}{\pi^3} \sum_{n=0}^{\infty} \frac{1}{(2n+1)^3} \cos \frac{(2n+1)\pi a}{2l}t \sin \frac{(2n+1)\pi}{2l}x
	\end{equation}
	
	\section{热方程}

	“计算题中最有可能考混合问题分离变量法”
	
	\subsection{傅里叶变换}
	\subsubsection{定义与反演公式}
	
	\paragraph{定义} 
	若 $f(x) \in L(-\infty, \infty)$，那么积分
	\begin{equation}
		\hat{f}(\lambda) = \frac{1}{\sqrt{2\pi}} \int_{-\infty}^{\infty} f(x) \mathrm{e}^{-\mathrm{i}\lambda x} \mathrm{d}x
	\end{equation}
	有意义，它定义了一个从 $f(x) \to \hat{f}(\lambda)$ 的变换。通常把上式左端积分所定义的变换称之为 $f(x)$ 的 \textbf{Fourier 变换}，记作 $\hat{f}$。$\hat{f}(\lambda)$ 也称为 $f(x)$ 的 Fourier 变式（或 Fourier 变换的像）。
	
	\paragraph{定理 (Fourier 积分定理)}
	若 $f(x) \in L(-\infty, \infty) \cap C^1(-\infty, \infty)$，那么我们有
	\begin{equation}
		f(x) = \lim_{N \to \infty} \frac{1}{\sqrt{2\pi}} \int_{-N}^{N} \hat{f}(\lambda) \mathrm{e}^{\mathrm{i}\lambda x} \mathrm{d}\lambda
	\end{equation}
	该公式称为 \textbf{反演公式}。左端的积分表示取 Cauchy 主值。
	
	\subsubsection{傅里叶变换的性质}
	
	可能考性质的证明与计算
	
	以下性质假设相关函数满足 Fourier 变换存在的条件。
	
	\begin{enumerate}
		% 1. 线性性质
		\item \textbf{线性性质} \\
		若 $f_i(x) \in L(-\infty, \infty), a_i \in \mathbb{C} \; (i=1,2)$，则
		\begin{equation}
			(a_1 f_1 + a_2 f_2)^{\wedge} = a_1 \hat{f}_1 + a_2 \hat{f}_2
		\end{equation}
		
		% 2. 微商性质
		\item \textbf{微商性质} \\
		若 $f(x), f'(x) \in C(-\infty, \infty) \cap L(-\infty, \infty)$，则
		\begin{equation}
			\left( \frac{\mathrm{d}f}{\mathrm{d}x} \right)^{\wedge} = \mathrm{i}\lambda \hat{f}
		\end{equation}
		
		% 3. 乘多项式
		\item \textbf{乘多项式} \\
		若 $f(x), xf(x) \in L(-\infty, \infty)$，则有
		\begin{equation}
			(xf(x))^{\wedge} = \mathrm{i} \frac{\mathrm{d}}{\mathrm{d}\lambda} \hat{f}(\lambda)
		\end{equation}
		
		% 4. 平移性质
		\item \textbf{平移性质} \\
		若 $f(x) \in L(-\infty, \infty)$，则
		\begin{equation}
			(f(x-a))^{\wedge}(\lambda) = \mathrm{e}^{-\mathrm{i}\lambda a} \hat{f}(\lambda)
		\end{equation}
		
		% 5. 伸缩性质
		\item \textbf{伸缩性质} \\
		若 $f(x) \in L(-\infty, \infty)$，则
		\begin{equation}
			(f(kx))^{\wedge}(\lambda) = \frac{1}{|k|} \hat{f}\left( \frac{\lambda}{k} \right), \quad k \neq 0
		\end{equation}
		
		% 6. 对称性质
		\item \textbf{对称性质} \\
		若 $f(x) \in L(-\infty, \infty)$，则
		\begin{equation}
			(f(x))^{\vee}(\lambda) = \hat{f}(-\lambda) = [f(-x)]^{\wedge}(\lambda)
		\end{equation}
		
		% 7. 卷积性质
		\item \textbf{卷积性质} \\
		若 $f(x), g(x) \in L(-\infty, \infty)$，定义卷积
		\[
		f * g(x) = \int_{-\infty}^{\infty} f(x-t)g(t) \mathrm{d}t \in L(-\infty, \infty)
		\]
		则有
		\begin{equation}
			(f * g)^{\wedge}(\lambda) = \sqrt{2\pi} \hat{f} \hat{g}
		\end{equation}
	\end{enumerate}
	
	\subsection{广义导数}
	
	可能会考一计算，还有若干定义
	\vspace{1cm}
	
	% 定义 1.3：基本空间与试验函数
	\noindent\textbf{定义 } \quad 如果 $\varphi(x) \in C_0^\infty(\mathbb{R})$，$\varphi_n(x) \in C_0^\infty(\mathbb{R})$ ($n=1,2,\cdots$)，且满足：
	\begin{enumerate}
		\item 存在 $M>0$，使得当 $|x| \geqslant M$ 时 $\varphi(x) \equiv 0$，$\varphi_n(x) \equiv 0$ ($n=1,2,\cdots$)；
		\item 对于任意非负整数 $k$，导数一致收敛，即
		\[
		\lim_{n \to \infty} \max_{x \in [-M, M]} \left| \varphi_n^{(k)}(x) - \varphi^{(k)}(x) \right| = 0, \quad k=0,1,2,\cdots
		\]
	\end{enumerate}
	则称序列 $\{\varphi_n\}$ 收敛于 $\varphi$。规定了上述收敛性的线性空间 $C_0^\infty(\mathbb{R})$ 称为 \textbf{基本空间} $\mathscr{D}(\mathbb{R})$。$\varphi \in \mathscr{D}(\mathbb{R})$ 称为 \textbf{试验函数}。
	
	\vspace{1em}
	
	% 定义 1.4：广义函数
	\noindent\textbf{定义 } \quad 如果 $f: \mathscr{D}(\mathbb{R}) \to \mathbb{R}$ 是线性连续泛函，则称 $f$ 是一个 \textbf{广义函数}。设 $\varphi \in \mathscr{D}(\mathbb{R})$ 是一个试验函数，用 $\langle f, \varphi \rangle$ 表示它所对应的数值，称为 \textbf{对偶积}。
	
	\vspace{1em}
	
	% 例 1：Delta 函数
	\noindent\textbf{例} \quad 使任意试验函数 $\varphi(x)$ 与它在 $x=0$ 处的函数值 $\varphi(0)$ 对应的广义函数称为 $\delta$ 函数，即 $\delta$ 函数满足
	\begin{equation}
		\langle \delta, \varphi \rangle = \varphi(0), \quad \forall \varphi \in \mathscr{D}(\mathbb{R})
	\end{equation}
	事实上，容易验证 $\varphi \to \varphi(0)$ 是 $\mathscr{D}(\mathbb{R})$ 上的线性连续泛函。这就是前面所说的 $\delta$ 函数的严格定义。
	
	\vspace{1em}
	
	% 定义 1.7：广义导数
	\noindent\textbf{定义} \quad 广义函数 $f$ 的微商 $f'$ 也是广义函数，用下式定义：
	\begin{equation}
		\langle f', \varphi \rangle = - \langle f, \varphi' \rangle, \quad \forall \varphi \in \mathscr{D}(\mathbb{R})
	\end{equation}
	按照定义 1.7，可以继续定义 $f'$ 的微商 $f''$，即
	\[
	\langle f'', \varphi \rangle = - \langle f', \varphi' \rangle = (-1)^2 \langle f, \varphi'' \rangle, \quad \forall \varphi \in \mathscr{D}(\mathbb{R})
	\]
	类似地，广义函数 $f$ 的 $k$ 阶微商 $f^{(k)}$ 定义为
	\begin{equation}
		\langle f^{(k)}, \varphi \rangle = (-1)^k \langle f, \varphi^{(k)} \rangle, \quad \forall \varphi \in \mathscr{D}(\mathbb{R}), \ k=1,2,\cdots
	\end{equation}
	大家知道，通常的函数不见得是可微的，但根据上述定义，每个广义函数都是无穷次可微的。
	
	\noindent 我们可以这样计算 $\delta$ 函数的弱导数：
	\[
	\langle \delta^{(k)}, \varphi \rangle = (-1)^k \langle \delta, \varphi^{(k)} \rangle = (-1)^k \varphi^{(k)}(0), \quad \forall \varphi \in \mathscr{D}(\mathbb{R}), \ k=1,2,\cdots
	\]
	再比如说，针对 Heaviside 函数 $H(x) = \begin{cases} 1 & , x \geqslant 0 \\ 0 & , x < 0 \end{cases}$ ，我们可以这样计算：
	\[
	\langle H', \varphi \rangle = - \langle H, \varphi' \rangle = - \int_{-\infty}^{+\infty} H(x)\varphi'(x)\mathrm{d}x = \varphi(0)\langle \delta, \varphi \rangle
	\]
	也就是说 $H' = \delta$
	
	\subsection{Poisson 公式}
	
	可能直接考察公式证明
	
	在热传导方程的初值问题中，我们通过对空间变量 $x$ 进行 Fourier 变换来求解。考虑如下柯西问题：
	\begin{equation}
		\begin{cases}
			\frac{\partial u}{\partial t} = a^2 \frac{\partial^2 u}{\partial x^2} + f(x,t) \\
			u(x,0) = \varphi(x)
		\end{cases}
	\end{equation}
	对方程两边关于变量 $x$ 作 Fourier 变换，利用 Fourier 变换的微商性质 $\widehat{u_{xx}} = (\mathrm{i}\lambda)^2 \hat{u} = -\lambda^2 \hat{u}$，我们将偏微分方程转化为关于时间 $t$ 的常微分方程初值问题：
	\begin{equation}
		\begin{cases}
			\frac{\mathrm{d}\hat{u}}{\mathrm{d}t} + a^2 \lambda^2 \hat{u} = \hat{f}(\lambda, t) \\
			\hat{u}|_{t=0} = \hat{\varphi}(\lambda)
		\end{cases}
	\end{equation}
	其中 $\hat{u}(\lambda, t)$ 为解 $u(x,t)$ 关于 $x$ 的 Fourier 变式。这是一个一阶线性非齐次常微分方程，利用常数变易法或积分因子法，可求得其解为：
	\begin{equation}
		\hat{u}(\lambda, t) = \hat{\varphi}(\lambda) \mathrm{e}^{-a^2 \lambda^2 t} + \int_0^t \hat{f}(\lambda, \tau) \mathrm{e}^{-a^2 \lambda^2 (t-\tau)} \mathrm{d}\tau
	\end{equation}
	为了得到原问题的解 $u(x,t)$，我们需要对上式两边进行 Fourier 逆变换。根据线性性质，我们有：
	\begin{equation}
		u(x,t) = \left( \hat{\varphi}(\lambda) \mathrm{e}^{-a^2 \lambda^2 t} \right)^{\vee} + \int_0^t \left( \hat{f}(\lambda, \tau) \mathrm{e}^{-a^2 \lambda^2 (t-\tau)} \right)^{\vee} \mathrm{d}\tau
	\end{equation}
	为了计算右端的逆变换，我们利用卷积定理：$(f * g)^{\wedge} = \sqrt{2\pi} \hat{f} \hat{g}$，即 $(\hat{f}\hat{g})^{\vee} = \frac{1}{\sqrt{2\pi}} f * g$。
	注意到高斯函数的变换对性质，若取
	\[
	g(x,t) = \frac{1}{a\sqrt{2t}} \mathrm{e}^{-\frac{x^2}{4a^2 t}}
	\]
	则其 Fourier 变换恰为 $\hat{g}(\lambda, t) = \mathrm{e}^{-a^2 \lambda^2 t}$。
	
	利用上述卷积关系，第一项的逆变换为：
	\begin{align*}
		\left( \hat{\varphi} \mathrm{e}^{-a^2 \lambda^2 t} \right)^{\vee} &= \frac{1}{\sqrt{2\pi}} \varphi * g \\
		&= \frac{1}{\sqrt{2\pi}} \int_{-\infty}^{\infty} \varphi(\xi) \frac{1}{a\sqrt{2t}} \mathrm{e}^{-\frac{(x-\xi)^2}{4a^2 t}} \mathrm{d}\xi \\
		&= \frac{1}{2a\sqrt{\pi t}} \int_{-\infty}^{\infty} \varphi(\xi) \mathrm{e}^{-\frac{(x-\xi)^2}{4a^2 t}} \mathrm{d}\xi
	\end{align*}
	同理，对于积分号内的第二项，利用时间平移后的性质，可得：
	\begin{align*}
		\left( \hat{f}(\lambda, \tau) \mathrm{e}^{-a^2 \lambda^2 (t-\tau)} \right)^{\vee} &= \frac{1}{2a\sqrt{\pi (t-\tau)}} \int_{-\infty}^{\infty} f(\xi, \tau) \mathrm{e}^{-\frac{(x-\xi)^2}{4a^2 (t-\tau)}} \mathrm{d}\xi
	\end{align*}
	将上述两部分结果代入 $u(x,t)$ 的表达式，整理即得热传导方程的解的积分表达式：
	\begin{equation}
		u(x,t) = \int_{-\infty}^{\infty} K(x-\xi, t) \varphi(\xi) \mathrm{d}\xi + \int_0^t \mathrm{d}\tau \int_{-\infty}^{\infty} K(x-\xi, t-\tau) f(\xi, \tau) \mathrm{d}\xi
	\end{equation}
	上式通常称为 \textbf{Poisson 公式}。其中函数 $K(x,t)$ 定义为：
	\begin{equation}
		K(x,t) = \begin{cases}
			\frac{1}{2a\sqrt{\pi t}} \mathrm{e}^{-\frac{x^2}{4a^2 t}}, & t > 0 \\
			0, & t \leqslant 0
		\end{cases}
	\end{equation}
	函数 $K(x,t)$（或记作 $\Gamma(x, t; \xi, \tau) = K(x-\xi, t-\tau)$）被称为热传导方程的 基本解 (Fundamental Solution) 或 热核 (Heat Kernel)。它描述了在 $t=0$ 时刻位于 $x=0$ 处的单位点热源在 $t>0$ 时刻产生的温度分布，且$K(x,t) \sim N(0, 2a^2t)$。
	
	\vspace{1em} % 增加一点垂直间距
	
	\noindent\textbf{定理（重要）}：
	若初值函数 $\varphi(x) \in C(-\infty, \infty)$ 且有界，源项 $f(x,t) \equiv 0$，则由 Poisson 公式确定的函数
	\begin{equation}
		u(x,t) = \int_{-\infty}^{\infty} K(x-\xi, t) \varphi(\xi) \mathrm{d}\xi = \frac{1}{2a\sqrt{\pi t}} \int_{-\infty}^{\infty} \mathrm{e}^{-\frac{(x-\xi)^2}{4a^2 t}} \varphi(\xi) \mathrm{d}\xi
	\end{equation}
	是热传导方程初值问题在 $t>0$ 时的有界解，且满足初始条件 $\lim_{t \to 0^+} u(x,t) = \varphi(x)$。
	
	\begin{proof}
		\textbf{第一步：验证 $u(x,t)$ 满足热传导方程 ($t>0$)}
		
		当 $t>0$ 时，由于热核 $K(x-\xi, t)$ 具有极好的光滑性（无穷次可微），且 $\varphi(\xi)$ 有界，我们可以直接在积分号下对参数 $x$ 和 $t$ 求导。
		注意到热核 $K$ 本身就是热传导方程的基本解，即满足：
		\[
		\frac{\partial K}{\partial t} - a^2 \frac{\partial^2 K}{\partial x^2} = 0
		\]
		因此：
		\begin{align*}
			\frac{\partial u}{\partial t} - a^2 \frac{\partial^2 u}{\partial x^2} &= \int_{-\infty}^{\infty} \left( \frac{\partial}{\partial t} - a^2 \frac{\partial^2}{\partial x^2} \right) K(x-\xi, t) \varphi(\xi) \mathrm{d}\xi \\
			&= \int_{-\infty}^{\infty} 0 \cdot \varphi(\xi) \mathrm{d}\xi = 0
		\end{align*}
		这说明 $u(x,t)$ 在 $t>0$ 时确实满足齐次热传导方程。
		
		\vspace{0.5em}
		
		\textbf{第二步：验证初始条件 ($t \to 0^+$)}
		
		我们需要证明对于任意 $x_0 \in (-\infty, \infty)$，有 $\lim_{x \to x_0, t \to 0^+} u(x,t) = \varphi(x_0)$。
		
		直接在原积分中取 $t \to 0$ 比较困难，因为 $K$ 会变成狄拉克 $\delta$ 函数。我们采用变量代换法，将 $t$ 从指数部分移到函数自变量中。
		令积分变量变换：
		\[
		\eta = \frac{\xi - x}{2a\sqrt{t}} \quad \Rightarrow \quad \xi = x + 2a\sqrt{t}\eta, \quad \mathrm{d}\xi = 2a\sqrt{t}\mathrm{d}\eta
		\]
		代入 $u(x,t)$ 的表达式，系数消去，积分上下限保持 $-\infty$ 到 $+\infty$：
		\begin{equation}
			u(x,t) = \frac{1}{\sqrt{\pi}} \int_{-\infty}^{\infty} \mathrm{e}^{-\eta^2} \varphi(x + 2a\sqrt{t}\eta) \mathrm{d}\eta
		\end{equation}
		现在我们可以对上式取极限。由于 $\varphi$ 有界，根据实变函数中的控制收敛定理（或利用一致收敛性），可以将极限运算放入积分号内部：
		\begin{align*}
			\lim_{x \to x_0, t \to 0^+} u(x,t) &= \frac{1}{\sqrt{\pi}} \int_{-\infty}^{\infty} \lim_{x \to x_0, t \to 0^+} \left[ \mathrm{e}^{-\eta^2} \varphi(x + 2a\sqrt{t}\eta) \right] \mathrm{d}\eta \\
			&= \frac{1}{\sqrt{\pi}} \int_{-\infty}^{\infty} \mathrm{e}^{-\eta^2} \varphi(x_0) \mathrm{d}\eta
		\end{align*}
		将常数 $\varphi(x_0)$ 提到积分号外，利用高斯积分公式 $\int_{-\infty}^{\infty} \mathrm{e}^{-\eta^2} \mathrm{d}\eta = \sqrt{\pi}$，我们得到：
		\[
		\lim_{x \to x_0, t \to 0^+} u(x,t) = \varphi(x_0) \cdot \frac{1}{\sqrt{\pi}} \cdot \sqrt{\pi} = \varphi(x_0)
		\]
		证毕。
	\end{proof}
	
	\subsection{基本解}
	
	2选1背诵
	
	在前文中我们导出了 Poisson 公式及其核函数 $K(x,t)$。在广义函数（分布）的理论框架下，热传导方程的基本解 $\Gamma(x,t; \xi, \tau)$ 有两种常见的等价定义方式。
	
	\subsubsection{定义一：基于非齐次方程（单位脉冲源）}
	
	设区域 $Q = \{(x,t) \mid -\infty < x < \infty, t > 0\}$。对于任意固定的点 $(\xi, \tau) \in Q$，如果函数 $u(x,t)$ 满足以下条件：
	\begin{enumerate}
		\item \textbf{正则性}：$u(x,t) \in L_{\text{loc}}(Q) \cap C(\bar{Q} \setminus \{(\xi, \tau)\})$，即在除去点 $(\xi, \tau)$ 之外的区域上连续；
		\item \textbf{方程与初值}：在广义函数的意义下满足
		\begin{equation*}
			\begin{cases}
				u_t - a^2 u_{xx} = \delta(x-\xi, t-\tau), & (x,t) \in Q \\
				u(x,0) = 0, & -\infty < x < \infty
			\end{cases}
		\end{equation*}
	\end{enumerate}
	则称 $u(x,t)$ 为热传导方程的\textbf{基本解}，记为 $\Gamma(x,t; \xi, \tau)$。
	\newline
	\textbf{注}：$L_{loc}$指局部可积函数空间，对于 $\Omega$ 内部的任意一个紧致子集（Compact Subset）$K$，函数 $f$ 在 $K$ 上都是可积的。
	
	$$\int_K |f(x)| \mathrm{d}x < \infty, \quad \forall \text{ 紧致集合 } K \subset \Omega$$

	\textbf{注}：
	
	$$\delta(x) = \begin{cases} 
		+\infty, & x = 0 \\
		0, & x \neq 0
	\end{cases}
	\quad \text{且} \quad \int_{-\infty}^{\infty} \delta(x) \mathrm{d}x = 1$$
	
	\subsubsection{定义二：基于齐次方程（单位初始脉冲）}
	
	设区域 $Q = \{(x,t) \mid -\infty < x < \infty, t > 0\}$。对于任意固定的 $\xi \in \mathbb{R}$，如果函数 $u(x,t)$ 满足：
	\begin{enumerate}
		\item \textbf{正则性}：$u(x,t) \in L_{\text{loc}}(Q)$；
		\item \textbf{方程与初值}：在广义函数的意义下满足
		\begin{equation*}
			\begin{cases}
				u_t - a^2 u_{xx} = 0, & (x,t) \in Q \\
				u \big|_{t=0} = \delta(x-\xi)
			\end{cases}
		\end{equation*}
	\end{enumerate}
	则称该解为热传导方程的\textbf{基本解}。
	
	\subsection{半无界问题}
	
	大概率不考，因为要结合poisson公式
	
	对于定义在半无界区域 $\bar{Q} = \{0 < x < \infty, 0 < t < \infty\}$ 上的热传导方程，我们需要结合边界条件来求解。常用的方法是延拓法（Method of Extension），即将半无界问题转化为全空间上的 Cauchy 问题求解。\textbf{看 $x=0$ 处给的是函数值还是导数值，判断使用奇延拓还是偶延拓}。
	
	\subsubsection{第一边值问题 (Dirichlet 问题)}
	
	考虑如下定解问题：
	\begin{equation}
		\begin{cases}
			u_t - a^2 u_{xx} = f(x,t), & 0 < x < \infty, t > 0 \\
			u(0,t) = g(t), & t > 0 \\
			u(x,0) = \varphi(x), & 0 \leqslant x < \infty
		\end{cases}
	\end{equation}
	
	\paragraph{1. 齐次边界情形 ($g(t) \equiv 0$)}
	当 $g(t) = 0$ 时，利用奇延拓方法。
	定义 $\varphi(x)$ 和 $f(x,t)$ 在全轴上的奇延拓函数 $\bar{\varphi}$ 和 $\bar{f}$：
	\[
	\bar{\varphi}(x) = \begin{cases} \varphi(x), & x \geqslant 0 \\ -\varphi(-x), & x < 0 \end{cases}, \quad
	\bar{f}(x,t) = \begin{cases} f(x,t), & x \geqslant 0, t \geqslant 0 \\ -f(-x,t), & x < 0, t \geqslant 0 \end{cases}
	\]
	显然 $\bar{\varphi}, \bar{f}$ 均为 $x$ 的奇函数。构造全空间上的 Cauchy 问题：
	\[
	\begin{cases}
		\bar{u}_t - a^2 \bar{u}_{xx} = \bar{f}(x,t), & -\infty < x < \infty, t > 0 \\
		\bar{u}(x,0) = \bar{\varphi}(x)
	\end{cases}
	\]
	由 Poisson 公式，其解为：
	\[
	\bar{u}(x,t) = \int_{-\infty}^{\infty} \Gamma(x-\xi, t) \bar{\varphi}(\xi) \mathrm{d}\xi + \int_0^t \mathrm{d}\tau \int_{-\infty}^{\infty} \Gamma(x-\xi, t-\tau) \bar{f}(\xi, \tau) \mathrm{d}\xi
	\]
	利用奇函数的性质，将积分区间 $(-\infty, \infty)$ 拆分为 $(-\infty, 0)$ 和 $(0, \infty)$。对于负半轴的积分，作变量代换 $\eta = -\xi$，利用 $\bar{\varphi}(-\eta) = -\varphi(\eta)$，可得：
	\begin{align*}
		\int_{-\infty}^{\infty} \Gamma(x-\xi, t) \bar{\varphi}(\xi) \mathrm{d}\xi 
		&= \int_{0}^{\infty} \Gamma(x-\xi, t) \varphi(\xi) \mathrm{d}\xi + \int_{-\infty}^{0} \Gamma(x-\xi, t) [-\varphi(-\xi)] \mathrm{d}\xi \\
		&= \int_{0}^{\infty} \left[ \Gamma(x-\xi, t) - \Gamma(x+\xi, t) \right] \varphi(\xi) \mathrm{d}\xi
	\end{align*}
	同理处理非齐次项。最终得到半无界区域上第一边值问题的解：
	\begin{equation}
		u(x,t) = \int_{0}^{\infty} G_D(x, \xi, t) \varphi(\xi) \mathrm{d}\xi + \int_0^t \mathrm{d}\tau \int_{0}^{\infty} G_D(x, \xi, t-\tau) f(\xi, \tau) \mathrm{d}\xi
	\end{equation}
	其中核函数 $G_D$（Dirichlet 格林函数）由两个基本解相减构成：
	\begin{equation}
		G_D(x, \xi, t) = \Gamma(x-\xi, t) - \Gamma(x+\xi, t)
	\end{equation}
	
	\paragraph{2. 非齐次边界情形 ($g(t) \not\equiv 0$)}
	作函数代换 $u(x,t) = v(x,t) + g(t)$。则 $v(x,t)$ 满足齐次边界条件的定解问题：
	\[
	\begin{cases}
		v_t - a^2 v_{xx} = f(x,t) - g'(t), & 0 < x < \infty \\
		v(0,t) = 0, & t > 0 \\
		v(x,0) = \varphi(x) - g(0), & 0 \leqslant x < \infty
	\end{cases}
	\]
	此时可直接套用上述齐次边界的解法求出 $v(x,t)$，进而得到原问题的解 $u(x,t)$。
	
	\subsubsection{第二边值问题 (Neumann 问题)}
	
	考虑边界条件为已知热流密度的情形：
	\begin{equation}
		\begin{cases}
			u_t - a^2 u_{xx} = f(x,t), & 0 < x < \infty, t > 0 \\
			u_x(0,t) = g(t), & t > 0 \\
			u(x,0) = \varphi(x), & 0 \leqslant x < \infty
		\end{cases}
	\end{equation}
	
	\paragraph{1. 绝热边界情形 ($g(t) \equiv 0$)}
	当 $u_x(0,t) = 0$ 时，利用偶延拓方法。
	定义 $\varphi, f$ 的偶延拓 $\bar{\varphi}, \bar{f}$（即 $\bar{\varphi}(-x) = \bar{\varphi}(x)$）。类似于第一边值问题的推导，由于偶函数性质，负半轴积分产生的项与正半轴相加。
	解的表达式为：
	\begin{equation}
		u(x,t) = \int_{0}^{\infty} G_N(x, \xi, t) \varphi(\xi) \mathrm{d}\xi + \int_0^t \mathrm{d}\tau \int_{0}^{\infty} G_N(x, \xi, t-\tau) f(\xi, \tau) \mathrm{d}\xi
	\end{equation}
	其中核函数 $G_N$（Neumann 格林函数）由两个基本解相加构成：
	\begin{equation}
		G_N(x, \xi, t) = \Gamma(x-\xi, t) + \Gamma(x+\xi, t)
	\end{equation}
	
	\paragraph{2. 非齐次边界情形 ($g(t) \not\equiv 0$)}
	作函数代换 $u(x,t) = v(x,t) + xg(t)$。
	验证边界条件：$u_x(0,t) = v_x(0,t) + g(t)$。若要 $u_x(0,t) = g(t)$，只需 $v_x(0,t) = 0$。
	此时 $v(x,t)$ 满足以下齐次 Neumann 问题：
	\[
	\begin{cases}
		v_t - a^2 v_{xx} = f(x,t) - xg'(t), & 0 < x < \infty \\
		v_x(0,t) = 0, & t > 0 \\
		v(x,0) = \varphi(x) - xg(0), & 0 \leqslant x < \infty
	\end{cases}
	\]
	求解 $v(x,t)$ 后即可得到原问题的解。
	
	\subsection{混合问题（分离变量法习题）}
	
	大概率考1题计算
	
	本节主要通过习题巩固有限区间 $0 < x < l$ 上热传导方程混合问题的求解。求解方法与波动方程类似，均采用分离变量法。
	
	\begin{enumerate}
		\item \textbf{齐次边界条件的求解} \\
		用分离变量法求解下列混合问题：
		
		\begin{enumerate}
			% 对应图片 9.(2)
			\item 
			\begin{equation}
				\begin{cases}
					u_t = a^2 u_{xx}, & 0 < x < \pi, \, t > 0 \\
					u|_{t=0} = \sin x, & 0 \leqslant x \leqslant \pi \\
					u_x|_{x=0} = u_x|_{x=\pi} = 0, & t > 0
				\end{cases}
			\end{equation}
			
			% 对应图片 9.(3)
			\item 
			\begin{equation}
				\begin{cases}
					u_t = a^2 u_{xx}, & 0 < x < l, \, t > 0 \\
					u|_{t=0} = x^2(l-x), & 0 \leqslant x \leqslant l \\
					u|_{x=0} = u|_{x=l} = 0, & t > 0
				\end{cases}
			\end{equation}
		\end{enumerate}
		
		\item \textbf{非齐次边界条件的求解} \\
		求解下列问题（提示：先通过函数代换 $u(x,t) = v(x,t) + w(x,t)$ 将边界条件齐次化）：
		
		\begin{enumerate}
			% 对应图片 9.(4)
			\item 
			\begin{equation}
				\begin{cases}
					u_t = a^2 u_{xx}, & 0 < x < l, \, t > 0 \\
					u|_{t=0} = 0, & 0 \leqslant x \leqslant l \\
					u|_{x=0} = 0, \quad u|_{x=l} = At, & t > 0
				\end{cases}
			\end{equation}
			
			% 对应图片 9.(6)
			\item 
			\begin{equation}
				\begin{cases}
					u_t - a^2 u_{xx} = 0, & 0 < x < l, \, t > 0 \\
					u|_{t=0} = 0, & 0 \leqslant x \leqslant l \\
					u_x|_{x=0} = 0, \quad u_x|_{x=l} = q, & t > 0
				\end{cases}
			\end{equation}
		\end{enumerate}
	\end{enumerate}
	
	\subsection{极值原理与最大模估计}
	
	极值原理相关与位势方程共同考20+内容
	
	对于抛物型方程，极值原理是研究解的唯一性、稳定性和有界性的重要工具。本节讨论热传导方程算子 $L u = \frac{\partial u}{\partial t} - a^2 \frac{\partial^2 u}{\partial x^2}$ 的极值性质。
	
	我们记 $Q = \{(x,t) \mid \text{空间域} \times (0, T]\}$ 为时空区域，$\bar{Q}$ 为其闭包。
	定义 抛物边界 $\Gamma$ 为区域 $Q$ 的下底面（$t=0$）与侧面边界的并集。
	
	\subsubsection{弱极值原理}
	
	\paragraph{定理  (弱极值原理)}
	设 $u \in C^{2,1}(Q) \cap C(\bar{Q})$ 且满足 $Lu = u_t - a^2 u_{xx} \leqslant 0$，则 $u$ 在 $\bar{Q}$ 上的最大值必在抛物边界 $\Gamma$ 上达到，即：
	\begin{equation}
		\max_{\bar{Q}} u(x,t) = \max_{\Gamma} u(x,t)
	\end{equation}
	
	\textbf{证明：}
	分两种情况讨论。
	
	\textbf{情形 1：严格不等式}
	假设 $Lu < 0$。我们断言 $u$ 不能在 $Q$ 内部（即 $t \in (0, T]$ 且 $x$ 在内部）达到最大值。
	反之，若在某点 $P_0(x_0, t_0) \in Q$ 取得最大值，则根据多元微积分极值存在的必要条件，在该点处有：
	\begin{enumerate}
		\item 关于 $x$ 的一阶导数为零：$\frac{\partial u}{\partial x}(P_0) = 0$；
		\item 关于 $x$ 的二阶导数非正：$\frac{\partial^2 u}{\partial x^2}(P_0) \leqslant 0$；
		\item 关于 $t$ 的导数（若是内部点 $t_0 < T$ 则为 0，若是上顶面 $t_0 = T$ 则非负）：$\frac{\partial u}{\partial t}(P_0) \geqslant 0$。
	\end{enumerate}
	将上述不等式代入算子 $L$：
	\[
	Lu(P_0) = \frac{\partial u}{\partial t}(P_0) - a^2 \frac{\partial^2 u}{\partial x^2}(P_0) \geqslant 0
	\]
	这与假设 $Lu < 0$ 矛盾。因此最大值只能在边界 $\Gamma$ 上取得。
	
	\textbf{情形 2：一般情况}
	对于 $Lu \leqslant 0$，我们采用微扰法。对任意 $\varepsilon > 0$，构造辅助函数：
	\[
	v(x,t) = u(x,t) - \varepsilon t
	\]
	计算其作用下的算子值：
	\[
	Lv = Lu - \varepsilon < -\varepsilon < 0
	\]
	根据情形 1 的结论， $v$ 的最大值必在 $\Gamma$ 上取得。即：
	\[
	\max_{\bar{Q}} v = \max_{\Gamma} v
	\]
	从而对于任意 $(x,t) \in \bar{Q}$：
	\[
	u(x,t) - \varepsilon t \leqslant \max_{\bar{Q}} v = \max_{\Gamma} v \leqslant \max_{\Gamma} u
	\]
	即 $u(x,t) \leqslant \max_{\Gamma} u + \varepsilon t \leqslant \max_{\Gamma} u + \varepsilon T$。
	令 $\varepsilon \to 0$，即得 $u(x,t) \leqslant \max_{\Gamma} u$。证毕。
	
	\paragraph{推论 1 (最小模原理)}
	若 $Lu \geqslant 0$，则 $u$ 在 $\bar{Q}$ 上的最小值必在 $\Gamma$ 上达到。
	特别地，若 $Lu = 0$（热传导方程的齐次解），则 $u$ 在 $\bar{Q}$ 上的最大值与最小值均在抛物边界 $\Gamma$ 上达到。
	（提示：令 $v = -u$，利用定理 3.1 即可）。
	
	\paragraph{推论 2 (比较原理)}
	设 $u, v \in C^{2,1}(Q) \cap C(\bar{Q})$，且满足
	\[
	\begin{cases}
		Lu \leqslant Lv, & (x,t) \in Q \\
		u \leqslant v, & (x,t) \in \Gamma
	\end{cases}
	\]
	则在整个 $\bar{Q}$ 上有 $u(x,t) \leqslant v(x,t)$。
	（证明：令 $w = u - v$，则 $Lw \leqslant 0$，应用定理 3.1 可知 $w$ 的最大值在 $\Gamma$ 上取得，而 $w|_\Gamma \leqslant 0$，故 $w \leqslant 0$）。
	
	\subsubsection{第一边值问题解的最大模估计}
	
	考虑有限区间上的第一边值问题：
	\begin{equation}
		\begin{cases}
			Lu = f(x,t), & (x,t) \in Q \\
			u|_{t=0} = \varphi(x), & 0 \leqslant x \leqslant l \\
			u|_{x=0} = g_1(t), \quad u|_{x=l} = g_2(t), & 0 \leqslant t \leqslant T
		\end{cases}
	\end{equation}
	
	\paragraph{定理 3.2}
	设 $u$ 是上述问题的解，则有如下估计：
	\begin{equation}
		\max_{\bar{Q}} |u(x,t)| \leqslant F T + B
	\end{equation}
	其中 $F = \sup_Q |f|$，$B = \max \{ \sup|\varphi|, \sup|g_1|, \sup|g_2| \}$（即边界上的最大值）。
	
	\textbf{证明：}
	构造辅助函数 $w_{\pm}(x,t) = (FT + B) \pm u(x,t)$。
	\begin{enumerate}
		\item \textbf{验证边界条件}：在 $\Gamma$ 上，由于 $|u| \leqslant B$，故
		\[
		w_{\pm}|_{\Gamma} = FT + B \pm u|_{\Gamma} \geqslant FT + B - B = FT \geqslant 0
		\]
		\item \textbf{验证方程}：
		\[
		L(w_{\pm}) = L(FT + B) \pm Lu = F \pm f(x,t) \geqslant F - |f| \geqslant 0
		\]
	\end{enumerate}
	根据推论 1（最小模原理），$Lw_{\pm} \geqslant 0$ 意味着 $w_{\pm}$ 的最小值在 $\Gamma$ 上取得。由于 $w_{\pm}|_{\Gamma} \geqslant 0$，故在整个 $\bar{Q}$ 上 $w_{\pm}(x,t) \geqslant 0$。
	即 $(FT + B) \pm u(x,t) \geqslant 0$，亦即 $|u(x,t)| \leqslant FT + B$。
	
	\textbf{注：} 此定理直接推出了第一边值问题解的唯一性和连续依赖于初边值条件的稳定性。
	
	\subsubsection{第二、三边值问题的最大模估计}
	
	第二边值问题（Neumann）和第三边值问题（Robin）可以统一写成如下形式的定解问题：
	\begin{equation}
		\begin{cases}
			L u = \frac{\partial u}{\partial t} - a^2 \frac{\partial^2 u}{\partial x^2} = f(x,t), & (x,t) \in Q \\
			u|_{t=0} = \varphi(x), & 0 \leqslant x \leqslant l \\
			\left[ -\frac{\partial u}{\partial x} + \alpha(t) u \right]_{x=0} = g_1(t), & 0 \leqslant t \leqslant T \\
			\left[ \frac{\partial u}{\partial x} + \beta(t) u \right]_{x=l} = g_2(t), & 0 \leqslant t \leqslant T
		\end{cases}
		\label{eq:mixed_bvp}
	\end{equation}
	其中假设 $\alpha(t) \geqslant 0, \beta(t) \geqslant 0$。当 $\alpha(t) \equiv 0$ 且 $\beta(t) \equiv 0$ 时，上述问题即为第二边值问题。
	
	为了得到解的最大模估计，我们需要先建立一个关键的引理（正性原理）。
	
	\paragraph{引理 3.3}
	设 $u(x,t) \in C^{2,1}(Q) \cap C^{1,0}(\bar{Q})$ 满足以下齐次不等式条件：
	\begin{equation}
		\begin{cases}
			L u \geqslant 0, & (x,t) \in Q \\
			u|_{t=0} \geqslant 0, & 0 \leqslant x \leqslant l \\
			\left[ -\frac{\partial u}{\partial x} + \alpha(t) u \right]_{x=0} \geqslant 0, & 0 \leqslant t \leqslant T \\
			\left[ \frac{\partial u}{\partial x} + \beta(t) u \right]_{x=l} \geqslant 0, & 0 \leqslant t \leqslant T
		\end{cases}
	\end{equation}
	则在整个闭区域 $\bar{Q}$ 上有 $u(x,t) \geqslant 0$。
	
	\textbf{证明：}
	我们采用反证法与微扰法结合。
	\begin{enumerate}
		\item \textbf{情形 1：严格不等式} \\
		假设边界条件满足严格不等式（$>0$）。若 $u$ 在 $\bar{Q}$ 上有负值，则其最小值必为负数。
		由定理 3.1（弱极值原理）知，$L u \geqslant 0$ 意味着最小值不能在 $Q$ 内部达到；又由 $u|_{t=0} \geqslant 0$，最小值也不能在底面达到。因此，最小值点 $P_0$ 只能位于侧边界 $x=0$ 或 $x=l$ 上。
		\begin{itemize}
			\item 若 $P_0$ 在 $x=0$ 上：在该点取最小值意味着 $\frac{\partial u}{\partial x} \geqslant 0$（向内导数非负）。若 $u(P_0) < 0$，考虑到 $\alpha(t) \geqslant 0$，则有：
			\[
			\left[ -\frac{\partial u}{\partial x} + \alpha(t) u \right]_{P_0} \leqslant 0 + \alpha(t) u(P_0) \leqslant 0
			\]
			这与假设的严格不等式 $>0$ 矛盾。
			\item 若 $P_0$ 在 $x=l$ 上：同理可导出矛盾。
		\end{itemize}
		因此 $u$ 不可能有负值，即 $u \geqslant 0$。
		
		\item \textbf{情形 2：一般情况} \\
		对于非严格不等式，构造辅助函数进行微扰。对任意 $\varepsilon > 0$，令
		\[
		v(x,t) = u(x,t) + \varepsilon \left[ 2a^2 t + \left( x - \frac{l}{2} \right)^2 \right]
		\]
		容易验证 $v$ 满足 $L v = L u + \varepsilon (2a^2 - a^2 \cdot 2) = L u \geqslant 0$。
		同时在边界上，由于微扰项及其导数的性质，可验证 $v$ 满足严格的边界不等式（例如在 $x=0$ 处，导数项产生的贡献与函数值项均使得结果 $>0$）。
		应用情形 1 的结论得 $v(x,t) \geqslant 0$。令 $\varepsilon \to 0$，即得 $u(x,t) \geqslant 0$。
	\end{enumerate}
	
	\paragraph{定理 3.4 (最大模估计)}
	设 $u(x,t)$ 是问题 (\ref{eq:mixed_bvp}) 的解，则有如下估计：
	\begin{equation}
		|u(x,t)| \leqslant C(F + B)
	\end{equation}
	其中 $F = \sup_Q |f|$，$B = \max \left\{ \sup|g_1|, \sup|g_2|, \sup|\varphi| \right\}$。
	常数 $C$ 仅依赖于 $a, l, T$，具体表达式为：
	\[
	C = \max \left\{ T, \, 1 + \frac{2a^2 T}{l} + \frac{l}{4} \right\}
	\]
	
	\textbf{证明：}
	我们需要构造一个“比较函数”来控制 $u$。引入辅助函数：
	\[
	w_{\pm}(x,t) = F t + B z(x,t) \pm u(x,t)
	\]
	其中 $z(x,t)$ 是为了处理导数边界条件而精心构造的抛物线型函数：
	\[
	z(x,t) = 1 + \frac{1}{l} \left[ 2a^2 t + \left( x - \frac{l}{2} \right)^2 \right]
	\]
	我们验证 $w_{\pm}$ 满足引理 3.3 的所有条件：
	\begin{enumerate}
		\item \textbf{方程算子}：
		\[
		L z = \frac{1}{l}(2a^2) - a^2 \frac{1}{l}(2) = 0 \implies L w_{\pm} = F \pm f \geqslant 0
		\]
		\item \textbf{初始条件} ($t=0$)：
		注意到 $z(x,0) \geqslant 1$，故 $w_{\pm}|_{t=0} = B z(x,0) \pm \varphi(x) \geqslant B - |\varphi| \geqslant 0$。
		\item \textbf{边界条件} ($x=0$)：
		计算 $z$ 在边界的性质：$z|_{x=0} \geqslant 1$，且 $-\frac{\partial z}{\partial x}|_{x=0} = -\frac{1}{l} 2(-\frac{l}{2}) = 1$。
		于是对于算子 $D_0 u = -\frac{\partial u}{\partial x} + \alpha(t)u$：
		\begin{align*}
			D_0 w_{\pm} &= F t \cdot \alpha(t) + B \left( -\frac{\partial z}{\partial x} + \alpha(t) z \right) \pm g_1(t) \\
			&\geqslant 0 + B(1 + \alpha(t) z) - B \\
			&= B \alpha(t) z \geqslant 0
		\end{align*}
		\item \textbf{边界条件} ($x=l$)：
		同理可证 $\left[ \frac{\partial w_{\pm}}{\partial x} + \beta(t) w_{\pm} \right]_{x=l} \geqslant 0$。
	\end{enumerate}
	由引理 3.3，在 $\bar{Q}$ 上 $w_{\pm}(x,t) \geqslant 0$，即 $|u(x,t)| \leqslant F t + B z(x,t)$。
	最后，放缩 $t \leqslant T$ 和 $z(x,t)$ 的最大值：
	\[
	\max_{\bar{Q}} z(x,t) = 1 + \frac{1}{l} \left( 2a^2 T + \frac{l^2}{4} \right) = 1 + \frac{2a^2 T}{l} + \frac{l}{4}
	\]
	取 $C$ 为 $t$ 的系数（即 $T$）与 $z$ 的上界中的较大者，即得证。
	
	\subsubsection{初值问题解的最大模估计（无界区域）}
	
	在带形区域 $Q = \{(x,t) \mid -\infty < x < \infty, 0 < t \leqslant T\}$ 上考虑初值问题：
	\begin{equation}
		\begin{cases}
			u_t - a^2 u_{xx} = f(x,t) \\
			u(x,0) = \varphi(x)
		\end{cases}
	\end{equation}
	
	\paragraph{定理 3.5}
	设 $u$ 是上述初值问题的有界解，则有估计：
	\begin{equation}
		\sup_Q |u(x,t)| \leqslant T \sup_Q |f(x,t)| + \sup_{-\infty < x < \infty} |\varphi(x)|
	\end{equation}
	
	\textbf{证明：}
	
	我们需要采用截断区域并构造增长主导函数的技巧。
	记 $F = \sup |f|, \Phi = \sup |\varphi|, K = \sup |u|$。
	考虑截断区域 $Q_L = \{(x,t) \mid |x| < L, 0 < t \leqslant T\}$。
	构造辅助函数：
	\[
	w(x,t) = FT + \Phi + \frac{K}{L^2}(x^2 + 2a^2t) \pm u(x,t)
	\]
	其中 $v_L(x,t) = \frac{K}{L^2}(x^2 + 2a^2t)$ 满足 $L(v_L) = \frac{K}{L^2}(2a^2) - a^2 \frac{K}{L^2}(2) = 0$。
	\begin{enumerate}
		\item \textbf{验证方程}：$Lw = L(FT+\Phi) + L(v_L) \pm Lu = F \pm f \geqslant 0$。
		\item \textbf{验证边界}：
		\begin{itemize}
			\item 在 $t=0$ 处：$w|_{t=0} = \Phi + \frac{K}{L^2}x^2 \pm \varphi(x) \geqslant 0$。
			\item 在 $x=\pm L$ 处：$v_L \geqslant \frac{K}{L^2}L^2 = K$。
			\[
			w|_{x=\pm L} = FT + \Phi + v_L \pm u \geqslant FT + \Phi + K - K \geqslant 0
			\]
		\end{itemize}
	\end{enumerate}
	由极值原理，在 $Q_L$ 上 $w(x,t) \geqslant 0$。
	固定任意点 $(x_0, t_0) \in Q$，当 $L$ 充分大时 $(x_0, t_0) \in Q_L$，此时有：
	\[
	|u(x_0, t_0)| \leqslant FT + \Phi + \frac{K}{L^2}(x_0^2 + 2a^2 t_0)
	\]
	令 $L \to \infty$，第三项趋于 0，得 $|u(x_0, t_0)| \leqslant FT + \Phi$。证毕。
	
	\textbf{注（Tychonoff 唯一性定理）：}
	对于无界区域上的初值问题，仅有上述有界解的唯一性是不够的。可以证明，只要解满足如下增长条件：
	\[
	|u(x,t)| \leqslant M e^{Nx^2} \quad (M, N \text{为常数})
	\]
	则解是唯一的。若不满足此条件，即使对于零初值的齐次方程，也存在非零解。
	
	\subsection{边值问题解的能量模估计}
	
	大概率考证明
	
	能量法不仅适用于抛物型方程，也适用于双曲型方程。它通过建立能量积分不等式，从而证明解的唯一性和稳定性。
	
	记 $Q_T = \{(x,t) \mid 0 < x < l, 0 < t \leqslant T\}$。在 $Q_T$ 上考虑混合问题：
	\begin{equation}
		\begin{cases}
			\frac{\partial u}{\partial t} - a^2 \frac{\partial^2 u}{\partial x^2} = f(x,t), & (x,t) \in Q_T \\
			u(x,0) = \varphi(x), & 0 \leqslant x \leqslant l \\
			u(0,t) = u(l,t) = 0, & 0 \leqslant t \leqslant T
		\end{cases}
		\label{eq:mixed_prob}
	\end{equation}
	
	\subsubsection{能量不等式}
	
	\textbf{定理}：
	设 $u \in C^{1,0}(\bar{Q}_T) \cap C^{2,1}(Q_T)$ 是问题 (\ref{eq:mixed_prob}) 的解，则有如下估计：
	\begin{equation}
		\sup_{0 \leqslant t \leqslant T} \int_0^l u^2(x,t) \mathrm{d}x + 2a^2 \int_0^T \int_0^l \left( \frac{\partial u}{\partial x} \right)^2 \mathrm{d}x \mathrm{d}t \leqslant M \left( \int_0^l \varphi^2(x) \mathrm{d}x + \int_0^T \int_0^l f^2(x,t) \mathrm{d}x \mathrm{d}t \right)
	\end{equation}
	其中常数 $M$ 仅与 $T$ 有关（具体为 $M = 2(1 + \mathrm{e}^T)$）。
	
	\subsubsection{证明过程}
	
	\paragraph{1. 构造能量积分}
	在方程两边同乘以 $u(x,t)$，并在区域 $Q_t = [0,l] \times [0,t]$ 上积分：
	\[
	\int_0^t \int_0^l u \frac{\partial u}{\partial \tau} \mathrm{d}x \mathrm{d}\tau - a^2 \int_0^t \int_0^l u \frac{\partial^2 u}{\partial x^2} \mathrm{d}x \mathrm{d}\tau = \int_0^t \int_0^l u f \mathrm{d}x \mathrm{d}\tau
	\]
	
	\paragraph{2. 分部积分与边界条件处理}
	\begin{itemize}
		\item \textbf{第一项}（关于时间的导数）：
		\[
		\int_0^t \int_0^l u \frac{\partial u}{\partial \tau} \mathrm{d}x \mathrm{d}\tau = \int_0^l \left( \int_0^t \frac{1}{2} \frac{\partial}{\partial \tau} (u^2) \mathrm{d}\tau \right) \mathrm{d}x = \frac{1}{2} \int_0^l u^2(x,t) \mathrm{d}x - \frac{1}{2} \int_0^l \varphi^2(x) \mathrm{d}x
		\]
		\item \textbf{第二项}（关于空间的导数）：
		利用分部积分公式：
		\[
		- \int_0^l u \frac{\partial^2 u}{\partial x^2} \mathrm{d}x = - \left[ u \frac{\partial u}{\partial x} \right]_{0}^{l} + \int_0^l \left( \frac{\partial u}{\partial x} \right)^2 \mathrm{d}x
		\]
		由于边界条件 $u(0,\tau) = u(l,\tau) = 0$，边界项为零。
	\end{itemize}
	将上述结果代回方程积分式，整理得：
	\begin{equation}
		\frac{1}{2} \int_0^l u^2(x,t) \mathrm{d}x + a^2 \int_0^t \int_0^l \left( \frac{\partial u}{\partial x} \right)^2 \mathrm{d}x \mathrm{d}\tau = \frac{1}{2} \int_0^l \varphi^2(x) \mathrm{d}x + \int_0^t \int_0^l u f \mathrm{d}x \mathrm{d}\tau
	\end{equation}
	
	\paragraph{3. 应用 Cauchy-Schwarz 不等式}
	对于右端的最后一项，利用不等式 $ab \leqslant \frac{1}{2}(a^2 + b^2)$：
	\[
	\int_0^t \int_0^l u f \mathrm{d}x \mathrm{d}\tau \leqslant \frac{1}{2} \int_0^t \int_0^l u^2 \mathrm{d}x \mathrm{d}\tau + \frac{1}{2} \int_0^t \int_0^l f^2 \mathrm{d}x \mathrm{d}\tau
	\]
	代入上式并移项整理，得到基本能量不等式：
	\begin{align}
		\frac{1}{2} \int_0^l u^2(x,t) \mathrm{d}x &+ a^2 \int_0^t \int_0^l \left( \frac{\partial u}{\partial x} \right)^2 \mathrm{d}x \mathrm{d}\tau \nonumber \\
		&\leqslant \frac{1}{2} \int_0^l \varphi^2(x) \mathrm{d}x + \frac{1}{2} \int_0^t \int_0^l f^2 \mathrm{d}x \mathrm{d}\tau + \frac{1}{2} \int_0^t \int_0^l u^2 \mathrm{d}x \mathrm{d}\tau
		\label{eq:energy_pre}
	\end{align}
	
	\paragraph{4. Gronwall 不等式的应用}
	为了处理不等式右端包含未知函数 $u^2$ 的积分项，我们引入辅助函数。记：
	\[
	\Omega(t) = \int_0^t \int_0^l u^2 \mathrm{d}x \mathrm{d}\tau, \quad F(t) = \int_0^l \varphi^2(x) \mathrm{d}x + \int_0^t \int_0^l f^2 \mathrm{d}x \mathrm{d}\tau
	\]
	显然 $\Omega(0)=0$，且 $\Omega'(t) = \int_0^l u^2(x,t) \mathrm{d}x$。
	忽略 (\ref{eq:energy_pre}) 左端非负的梯度项 $a^2 \int (\dots)^2$，我们有：
	\[
	\frac{1}{2} \Omega'(t) \leqslant \frac{1}{2} F(t) + \frac{1}{2} \Omega(t) \implies \frac{\mathrm{d}\Omega}{\mathrm{d}t} \leqslant F(t) + \Omega(t)
	\]
	由 Gronwall 不等式（微分形式），可得：
	\[
	\Omega(t) \leqslant \mathrm{e}^t \int_0^t \mathrm{e}^{-\tau} F(\tau) \mathrm{d}\tau + \Omega(0) \mathrm{e}^t
	\]
	考虑到 $F(t)$ 是 $t$ 的增函数，简单的放缩（如书中所述）可得更宽松的界：$\Omega(t) \leqslant \mathrm{e}^t F(t)$。
	将此界代回 (\ref{eq:energy_pre}) 式：
	\begin{align*}
		\frac{1}{2} \int_0^l u^2 \mathrm{d}x + a^2 \int_0^t \int_0^l \left( \frac{\partial u}{\partial x} \right)^2 \mathrm{d}x \mathrm{d}\tau 
		&\leqslant \frac{1}{2} F(t) + \frac{1}{2} \Omega(t) \\
		&\leqslant \frac{1}{2} F(T) + \frac{1}{2} \mathrm{e}^T F(T) \\
		&= \frac{1}{2} (1 + \mathrm{e}^T) F(T)
	\end{align*}
	最后两边乘以 2，并取上确界 $\sup_{0 \leqslant t \leqslant T}$，即得定理结论。
	
	\section{位势方程}
	
	证明基本解是基本解重要
	
	已知格林函数的时候如何求解位势方程
	
	Dirichlet问题的解，2维3维选一个考
	
	
	
	\subsection{基本解与格林函数}
	
	Laplace 方程 $\Delta u = 0$ 和 Poisson 方程 $-\Delta u = f$ 是位势理论的核心。求解这类方程的关键工具是基本解和格林公式。
	
	\subsubsection{Laplace 方程的基本解}
	
	\paragraph{定义 :基本解}
	如果函数 $U \in L_{\text{loc}}(\mathbb{R}^n)$ 在广义函数意义下满足
	\begin{equation}
		-\Delta U = \delta(x-\xi), \quad x, \xi \in \mathbb{R}^n
	\end{equation}
	即对于任意测试函数 $\varphi \in \mathscr{D}(\mathbb{R}^n)$，有
	\[
	\int_{\mathbb{R}^n} U(x) (-\Delta \varphi(x)) \mathrm{d}x = \varphi(\xi)
	\]
	则称 $U(x)$（或记为 $\Gamma(x, \xi)$）为 Laplace 方程的\textbf{基本解}。
	
	\paragraph{基本解的显式表达式}
	对于 $n$ 维 Laplace 方程，其基本解 $\Gamma(x, \xi) = \Gamma(|x-\xi|)$ 仅依赖于两点间的距离 $r = |x-\xi|$。具体形式如下：
	
	\begin{enumerate}
		\item \textbf{三维情形 ($n=3$)}：
		\begin{equation}
			\Gamma(x, \xi) = \frac{1}{4\pi r} = \frac{1}{4\pi |x-\xi|}
		\end{equation}
		
		\item \textbf{二维情形 ($n=2$)}：
		\begin{equation}
			\Gamma(x, \xi) = \frac{1}{2\pi} \ln \frac{1}{r} = \frac{1}{2\pi} \ln \frac{1}{|x-\xi|}
		\end{equation}
		
		\item \textbf{一般情形 ($n > 2$)}：
		\begin{equation}
			\Gamma(x, \xi) = \frac{1}{(n-2)\omega_n} \cdot \frac{1}{|x-\xi|^{n-2}}
		\end{equation}
		其中 $\omega_n = \frac{2\pi^{n/2}}{\Gamma(n/2)}$ 是 $n$ 维空间中单位球面的面积（例如 $\omega_3 = 4\pi$）。
	\end{enumerate}
	
	\subsubsection{格林公式及基本解的验证}
	
	设 $\Omega \subset \mathbb{R}^n$ 是有界区域，边界 $\partial \Omega$ 分段光滑。$n$ 是 $\partial \Omega$ 的单位外法向量。
	
	\paragraph{定理: 格林第二公式}
	设 $u, v \in C^2(\Omega) \cap C^1(\bar{\Omega})$，则有
	\begin{equation}
		\int_{\Omega} (u \Delta v - v \Delta u) \mathrm{d}x = \int_{\partial \Omega} \left( u \frac{\partial v}{\partial n} - v \frac{\partial u}{\partial n} \right) \mathrm{d}S
	\end{equation}
	该公式反映了区域内部的拉普拉斯算子与边界上的法向导数之间的对偶关系。
	
	\paragraph{定理:}由 4.1.1定义的函数 $\Gamma(x;\xi)$ 是 Laplace 方程的基本解。\textbf{重要}
	
	\begin{proof}
		由于 $\Gamma(x,y;\xi,\eta)$ 在 $(x,y) = (\xi,\eta)$ 时有奇性，故考虑区域
		\[
		\Omega_\varepsilon
		= \left\{ (x,y)\in\mathbb{R}^2
		\;\middle|\;
		\varepsilon^2 < (x-\xi)^2 + (y-\eta)^2 < \frac{1}{\varepsilon^2}
		\right\}.
		\]
		对于任意 $\varphi \in \mathcal{D}(\mathbb{R}^2)$，有
		\[
		\iint_{\mathbb{R}^2} \Gamma(x,y;\xi,\eta)\,(-\Delta\varphi(x,y))\,dxdy
		= \lim_{\varepsilon\to0^+}
		\iint_{\Omega_\varepsilon}
		\Gamma(x,y;\xi,\eta)\,(-\Delta\varphi(x,y))\,dxdy.
		\]
		
		在 $\Omega_\varepsilon$ 上应用 Green 公式，其中取
		\[
		u(x,y) = \varphi(x,y), \qquad
		v(x,y) = \Gamma(x,y;\xi,\eta),
		\]
		并注意到在 $\Omega_\varepsilon$ 上
		\[
		\Delta \Gamma(x,y;\xi,\eta) = 0.
		\]
		由 $\varphi \in \mathcal{D}(\mathbb{R}^2)$ 可知，当 $\varepsilon$ 充分小时，
		$\varphi(x,y)$ 在 $(x-\xi)^2+(y-\eta)^2 = \frac{1}{\varepsilon^2}$ 附近恒为零，
		因而有
		\begin{equation}
			\iint_{\mathbb{R}^2} \Gamma(x,y;\xi,\eta)\,(-\Delta\varphi(x,y))\,dxdy
			=
			\lim_{\varepsilon\to0^+}
			\int_{\rho=\varepsilon}
			\left(
			\Gamma(x,y;\xi,\eta)\frac{\partial\varphi(x,y)}{\partial\rho}
			-
			\frac{\partial\Gamma(x,y;\xi,\eta)}{\partial\rho}\,\varphi(x,y)
			\right) dl,
			\tag{1.10}
		\end{equation}
		其中
		\[
		\rho = \sqrt{(x-\xi)^2 + (y-\eta)^2}.
		\]
		
		利用 $\Gamma(x,y;\xi,\eta)$ 的表达式 (1.8)，有
		\[
		\left|
		\int_{\rho=\varepsilon}
		\Gamma(x,y;\xi,\eta)\frac{\partial\varphi(x,y)}{\partial\rho}\,dl
		\right|
		\le
		\int_{\rho=\varepsilon}
		\frac{1}{2\pi}
		\left|\frac{\partial\varphi(x,y)}{\partial\rho}\right|
		\ln\frac{1}{\rho}\,dl
		\]
		\[
		\le
		\max|\nabla\varphi|\cdot \varepsilon\ln\frac{1}{\varepsilon}
		\longrightarrow 0,
		\qquad \text{当 } \varepsilon\to0^+.
		\]
		
		另一方面，
		\[
		-
		\int_{\rho=\varepsilon}
		\frac{\partial\Gamma(x,y;\xi,\eta)}{\partial\rho}\,
		\varphi(x,y)\,dl
		=
		\frac{1}{2\pi}
		\int_{\rho=\varepsilon}
		\frac{1}{\rho}\,\varphi(x,y)\,dl
		\]
		\[
		=
		\frac{1}{2\pi\varepsilon}
		\int_{\rho=\varepsilon}
		\varphi(x,y)\,dl
		\longrightarrow
		\varphi(\xi,\eta),
		\qquad \text{当 } \varepsilon\to0^+.
		\]
		
		代入 (1.10) 式，得到
		\[
		\iint_{\mathbb{R}^2}
		\Gamma(x,y;\xi,\eta)\,(-\Delta\varphi(x,y))\,dxdy
		=
		\varphi(\xi,\eta).
		\]
		
		这就是所要证的。
	\end{proof}
	
	\subsubsection{解的积分表示 (积分表示定理)}
	
	利用基本解和格林公式，我们可以给出 Poisson 方程解的积分表达形式。这被称为格林表示公式或第三格林公式。
	
	\paragraph{定理}
	设 $u \in C^2(\Omega) \cap C^1(\bar{\Omega})$，则对于任意一点 $\xi \in \Omega$，我们有：
	\begin{equation}
		u(\xi) = - \int_{\Omega} \Gamma(x, \xi) \Delta u(x) \mathrm{d}x + \int_{\partial \Omega} \left[ \Gamma(x, \xi) \frac{\partial u(x)}{\partial n} - u(x) \frac{\partial \Gamma(x, \xi)}{\partial n_x} \right] \mathrm{d}S_x
	\end{equation}
	
	\subsubsection{格林函数 (Green Function)}
	
	为了消除积分表示定理中未知的边界法向导数项 $\frac{\partial u}{\partial n}$，我们引入格林函数的概念。格林函数可以看作是针对特定区域 $\Omega$ 及特定边界条件修正后的基本解。
	
	\paragraph{定义  : 格林函数}
	定义于 $\bar{\Omega} \times \bar{\Omega} \setminus \{ (x,y)=(\xi,\eta) \}$ 的连续函数
	\begin{equation}
		G(x,y; \xi,\eta) = \Gamma(x,y; \xi,\eta) + g(x,y; \xi,\eta)
	\end{equation}
	称为 Dirichlet 问题的 \textbf{Green 函数}，如果其中的 $g(x,y; \xi,\eta)$ 在 $\Omega \times \Omega$ 上二次连续可微，而且 $G(x,y; \xi,\eta)$ 满足以下问题（在广义函数意义下）：
	\begin{equation}
		\begin{cases}
			-\Delta_{(x,y)} G(x,y; \xi,\eta) = \delta(x-\xi, y-\eta), & (x,y) \in \Omega \\
			G(x,y; \xi,\eta) \big|_{\partial \Omega} = 0. &
		\end{cases}
	\end{equation}
	
	由上述定义可知，修正项（正则部分）$g(x,y; \xi,\eta)$ 满足以下边值问题：
	\begin{equation}
		\begin{cases}
			-\Delta_{(x,y)} g = 0, & (x,y) \in \Omega \\
			g \big|_{\partial \Omega} = -\Gamma(x,y; \xi,\eta) \big|_{\partial \Omega}. &
		\end{cases}
	\end{equation}
	这表明 $g$ 是 $\Omega$ 内的调和函数，其作用是抵消基本解 $\Gamma$ 在边界上的值，使得 $G$ 在边界上为零。
	
	\paragraph{格林函数的性质}
	下面给出 Green 函数的一些重要性质：
	
	\begin{enumerate}
		\item[$1^\circ$] \textbf{奇点性质与调和性}：
		当 $(x,y) \in \Omega \setminus \{ (\xi,\eta) \}$ 时，
		\[
		\Delta_{(x,y)} G(x,y; \xi,\eta) = 0.
		\]
		当 $(x,y) \to (\xi,\eta)$ 时，Green 函数 $G(x,y; \xi,\eta)$ 与基本解 $\Gamma(x,y; \xi,\eta)$ 是等价的。在二维情形下，即与
		\[
		\frac{1}{2\pi} \ln \left[ (x-\xi)^2 + (y-\eta)^2 \right]^{-1}
		\]
		等价。这说明格林函数保留了基本解在源点处的奇性结构。
		
		\item[$2^\circ$] \textbf{正性与上界估计}：
		当 $(x,y) \in \Omega \setminus \{ (\xi,\eta) \}$ 时，有
		\begin{equation}
			0 < G(x,y; \xi,\eta) \le \frac{1}{2\pi} \ln \frac{d}{\sqrt{(x-\xi)^2 + (y-\eta)^2}}
		\end{equation}
		其中 $d$ 是区域 $\Omega$ 的直径。
		
		\textbf{注}：这个性质的证明（特别是 $G>0$）需要利用\textbf{极值原理}。物理上，这对应于在 $\xi$ 处放置正点电荷，接地导体壳内部电势恒为正。
		
		\item[$3^\circ$] \textbf{对称性 (互惠性)}：
		Green 函数关于源点和场点具有对称性，即
		\begin{equation}
			G(x,y; \xi,\eta) = G(\xi,\eta; x,y).
		\end{equation}
		这意味着源点和观测点互换，产生的物理效应（如电势）不变。
		
		\item[$4^\circ$] \textbf{边界法向导数积分}：
		\begin{equation}
			\int_{\partial \Omega} \frac{\partial G(x,y; \xi,\eta)}{\partial n} \mathrm{d}l = -1.
		\end{equation}
		该性质可由方程 $-\Delta G = \delta$ 在 $\Omega$ 上积分并利用高斯公式得到（$-\int_{\partial \Omega} \frac{\partial G}{\partial n} dl = \iint_\Omega \delta d\sigma = 1$）。
	\end{enumerate}
	
	\paragraph{注记 : 利用 Green 函数求解 Dirichlet 问题}
	
	引入 Green 函数的主要动机是求解 Dirichlet 边值问题。已知区域 $\Omega$ 上的 Green 函数 $G(x, \xi)$ 后，对于如下 Dirichlet 问题：
	\begin{equation}
		\begin{cases}
			-\Delta u = f, & x \in \Omega \\
			u = \varphi, & x \in \partial \Omega
		\end{cases}
	\end{equation}
	其解 $u(\xi)$ 可直接由以下积分公式给出：
	\begin{equation}
		u(\xi) = \int_{\Omega} G(x, \xi) f(x) \mathrm{d}x - \int_{\partial \Omega} \varphi(x) \frac{\partial G(x, \xi)}{\partial n} \mathrm{d}S_x, \quad \xi \in \Omega
	\end{equation}
	
	\textbf{公式推导说明：}
	回顾解的积分表示（第三 Green 公式）：
	\[
	u(\xi) = - \int_{\Omega} \Gamma \Delta u \mathrm{d}x + \int_{\partial \Omega} \left( \Gamma \frac{\partial u}{\partial n} - u \frac{\partial \Gamma}{\partial n} \right) \mathrm{d}S
	\]
	我们将 $G(x, \xi) = \Gamma(x, \xi) + g(x, \xi)$ 代入 Green 第二公式，并利用 $g$ 在 $\Omega$ 内调和 ($\Delta g = 0$) 的性质，可以推导出类似的表示。
	关键在于，由于 Green 函数在边界上满足 $G|_{\partial \Omega} = 0$，使得上式中包含未知量 $\frac{\partial u}{\partial n}$ 的项 $\int_{\partial \Omega} G \frac{\partial u}{\partial n} \mathrm{d}S$ 自动消失，从而使解完全由已知数据 ($f$ 和 $\varphi$) 确定。
	
	\begin{itemize}
		\item \textbf{特例 (Laplace 方程)}：若 $f \equiv 0$，则解仅由边界积分给出：
		\begin{equation}
			u(\xi) = - \int_{\partial \Omega} \varphi(x) \frac{\partial G(x, \xi)}{\partial n} \mathrm{d}S_x
		\end{equation}
		这即是广义的 \textbf{Poisson 积分公式}。
	\end{itemize}


	\subsection{二维三维的Green函数及Poisson公式}
	
	\subsubsection{二维情形：圆域上的 Green 函数与 Poisson 公式}
	
	设 $\Omega = B_a(0) \subset \mathbb{R}^2$ 是以原点为心、半径为 $a$ 的圆盘。
	采用极坐标系：设源点 $\xi$ 的坐标为 $(\rho, \alpha)$，场点 $x$ 的坐标为 $(r, \theta)$。
	
	\paragraph{1. Green 函数的构造}
	利用镜像法，源点 $P(\rho, \alpha)$ 关于圆周的镜像点 $P^*$ 的极坐标为 $(\rho^*, \alpha)$，满足 $\rho \rho^* = a^2$，即 $\rho^* = \frac{a^2}{\rho}$。
	
	二维 Laplace 方程的基本解为 $\frac{1}{2\pi} \ln \frac{1}{r}$。圆域 Green 函数设为：
	\begin{equation}
		G(r, \theta; \rho, \alpha) = \frac{1}{2\pi} \ln \frac{1}{|x-\xi|} - \frac{1}{2\pi} \ln \frac{C}{|x-\xi^*|}.
	\end{equation}
	为使边界上为 0，当 $r=a$ 时，利用相似三角形性质 $|x-\xi| = \frac{\rho}{a} |x-\xi^*|$，我们需要常数 $C = \frac{a}{\rho}$。
	
	利用余弦定理展开距离项 $|x-\xi|^2 = r^2 + \rho^2 - 2r\rho \cos(\theta-\alpha)$，Green 函数可写为对数形式：
	\begin{equation}
		G(r, \theta; \rho, \alpha) = \frac{1}{4\pi} \ln \frac{1}{r^2 + \rho^2 - 2r\rho \cos(\theta-\alpha)} - \frac{1}{4\pi} \ln \frac{a^2/\rho^2}{r^2 + (\frac{a^2}{\rho})^2 - 2r \frac{a^2}{\rho} \cos(\theta-\alpha)}.
	\end{equation}
	Green 函数的常用显式为：
	\begin{equation}
		G(r, \theta; \rho, \alpha) = \frac{1}{4\pi} \ln \frac{\rho^2 \left[ r^2 + (\frac{a^2}{\rho})^2 - 2r \frac{a^2}{\rho} \cos(\theta-\alpha) \right]}{a^2 \left[ r^2 + \rho^2 - 2r\rho \cos(\theta-\alpha) \right]}.
	\end{equation}
	或者更为紧凑的形式（用于求导）：
	\[
	G = \frac{1}{4\pi} \ln \frac{\rho^2 r^2 + a^4 - 2r\rho a^2 \cos(\theta-\alpha)}{a^2 (r^2 + \rho^2 - 2r\rho \cos(\theta-\alpha))}.
	\]
	
	\paragraph{2. 边界法向导数的计算}
	在圆周 $\partial B_a$ 上，外法向方向即为径向 $r$。我们需要计算 $\frac{\partial G}{\partial n} \big|_{r=a} = \frac{\partial G}{\partial r} \big|_{r=a}$。
	
	记 $A = \rho^2 r^2 + a^4 - 2r\rho a^2 \cos \gamma$（分子项相关），$B = r^2 + \rho^2 - 2r\rho \cos \gamma$（分母项相关），其中 $\gamma = \theta - \alpha$。
	则 $G = \frac{1}{4\pi} (\ln A - \ln B - \ln a^2)$。
	
	对 $r$ 求偏导：
	\[
	\frac{\partial G}{\partial r} = \frac{1}{4\pi} \left( \frac{A'}{A} - \frac{B'}{B} \right) 
	= \frac{1}{4\pi} \left( \frac{2r\rho^2 - 2\rho a^2 \cos\gamma}{\rho^2 r^2 + a^4 - 2r\rho a^2 \cos\gamma} - \frac{2r - 2\rho \cos\gamma}{r^2 + \rho^2 - 2r\rho \cos\gamma} \right).
	\]
	
	当 $r=a$ 时，注意到此时分子 $A$ 和分母 $B$ 存在倍数关系：
	\[
	A|_{r=a} = \rho^2 a^2 + a^4 - 2a^3 \rho \cos\gamma = a^2 (\rho^2 + a^2 - 2a\rho \cos\gamma) = a^2 B|_{r=a}.
	\]
	代入导数表达式：
	\[
	\begin{aligned}
		\frac{\partial G}{\partial r} \bigg|_{r=a} &= \frac{1}{4\pi} \left( \frac{2a\rho^2 - 2\rho a^2 \cos\gamma}{a^2 B|_{r=a}} - \frac{2a - 2\rho \cos\gamma}{B|_{r=a}} \right) \\
		&= \frac{1}{4\pi B|_{r=a}} \left( \frac{2a\rho^2 - 2\rho a^2 \cos\gamma}{a^2} - (2a - 2\rho \cos\gamma) \right) \\
		&= \frac{1}{4\pi B|_{r=a}} \left( \frac{2\rho^2}{a} - 2\rho \cos\gamma - 2a + 2\rho \cos\gamma \right) \\
		&= \frac{1}{4\pi (a^2 + \rho^2 - 2a\rho \cos\gamma)} \left( \frac{2\rho^2 - 2a^2}{a} \right) \\
		&= \frac{1}{2\pi a} \frac{\rho^2 - a^2}{a^2 + \rho^2 - 2a\rho \cos(\theta - \alpha)}.
	\end{aligned}
	\]
	
	考虑圆上的 Dirichlet 问题：
	\begin{equation}
		\begin{cases}
			-\Delta u = f(x,y), & (x,y) \in B_a \\
			u \big|_{\partial B_a} = \varphi.
		\end{cases} 
	\end{equation}
	
	利用 Green 公式 $u = \int_{\Omega} G f \mathrm{d}x - \int_{\partial \Omega} \varphi \frac{\partial G}{\partial n} \mathrm{d}S$，
	我们得到解 $u(\rho, \theta)$ 的完整表达式：
	
	\begin{equation}
		\begin{aligned}
			u(\rho, \theta) = & \frac{1}{4\pi} \int_0^a r \mathrm{d}r \int_0^{2\pi} \ln \frac{\rho^2 r^2 + a^4 - 2r\rho a^2 \cos(\theta-\alpha)}{a^2 [ r^2 + \rho^2 - 2r\rho \cos(\theta-\alpha) ]} f(r, \alpha) \mathrm{d}\alpha \\
			& + \frac{1}{2\pi} \int_0^{2\pi} \frac{a^2 - \rho^2}{a^2 + \rho^2 - 2a\rho \cos(\theta - \alpha)} \varphi(\alpha) \mathrm{d}\alpha.
		\end{aligned} 
	\end{equation}
	
	\paragraph{3. 圆上的 Poisson 公式}
	Dirichlet 问题解的积分表示为：
	\[
	u(\rho, \alpha) = - \int_{\partial B_a} u(a, \theta) \frac{\partial G}{\partial n} \mathrm{d}S.
	\]
	在极坐标下，$\mathrm{d}S = a \mathrm{d}\theta$。将上述导数代入（注意负号使得 $\rho^2 - a^2$ 变为 $a^2 - \rho^2$）：
	\begin{equation}
		u(\rho, \alpha) = - \int_{0}^{2\pi} \varphi(\theta) \left[ \frac{1}{2\pi a} \frac{\rho^2 - a^2}{a^2 + \rho^2 - 2a\rho \cos(\theta - \alpha)} \right] a \mathrm{d}\theta.
	\end{equation}
	消去 $a$，整理得\textbf{圆上的 Poisson 公式}：
	\begin{equation}
		u(\rho, \alpha) = \frac{1}{2\pi} \int_{0}^{2\pi} \frac{a^2 - \rho^2}{a^2 + \rho^2 - 2a\rho \cos(\theta - \alpha)} \varphi(\theta) \mathrm{d}\theta.
	\end{equation}
	
	\begin{note}
		其中核函数 $P_r(\theta-\alpha) = \frac{1}{2\pi} \frac{a^2 - \rho^2}{a^2 + \rho^2 - 2a\rho \cos(\theta - \alpha)}$ 称为 Poisson 核。
		性质：
		\begin{enumerate}
			\item $P_r > 0$ （因为 $\rho < a$）；
			\item $\int_{0}^{2\pi} P_r(\phi) \mathrm{d}\phi = 1$；
			\item 当 $\rho \to a$ 时，$P_r$ 趋向于 $\delta$ 函数，使得 $u(\rho, \alpha)$ 逼近边界值 $\varphi(\alpha)$。
		\end{enumerate}
	\end{note}
	
	\subsubsection{三维情形：球域上的 Green 函数与 Poisson 公式}
	
	设 $\Omega = B_a(0) \subset \mathbb{R}^3$ 是半径为 $a$ 的球体。
	
	\paragraph{1. Green 函数的构造}
	利用镜像法，取 $\xi \in B_a(0)$ 为源点（$\xi \neq 0$），其关于球面的镜像点为 $\xi^* = \frac{a^2}{|\xi|^2} \xi$。
	球域上的 Green 函数由基本解和镜像产生的修正项组成：
	\begin{equation}
		G(x; \xi) = \frac{1}{4\pi} \left( \frac{1}{|x-\xi|} - \frac{a}{|\xi|} \frac{1}{|x-\xi^*|} \right).
	\end{equation}
	记 $r = |x-\xi|$，$r^* = |x-\xi^*|$，$\rho = |\xi|$。上式可简写为：
	\[
	G(x; \xi) = \frac{1}{4\pi} \left( \frac{1}{r} - \frac{a}{\rho} \frac{1}{r^*} \right).
	\]
	已验证当 $x \in \partial B_a$ 时，由相似关系 $r^* = \frac{a}{\rho} r$ 可得 $G(x; \xi) = 0$。
	
	\paragraph{2. 边界法向导数的计算（推导过程）}
	为了得到解的积分表示，我们需要计算 $G$ 在边界上的外法向导数 $\frac{\partial G}{\partial n} \big|_{|x|=a}$。
	
	注意外法向量 $n = \frac{x}{a}$（因为球心在原点）。根据链式法则，法向导数为：
	\[
	\frac{\partial G}{\partial n} = \nabla_x G \cdot n = \nabla_x G \cdot \frac{x}{a}.
	\]
	
	\textbf{第一步：计算梯度 $\nabla_x G$}
	利用公式 $\nabla_x (\frac{1}{r}) = -\frac{x-\xi}{r^3}$，有：
	\[
	\nabla_x G = \frac{1}{4\pi} \left[ -\frac{x-\xi}{r^3} - \frac{a}{\rho} \left( -\frac{x-\xi^*}{(r^*)^3} \right) \right] 
	= \frac{1}{4\pi} \left[ \frac{\xi-x}{r^3} + \frac{a}{\rho} \frac{x-\xi^*}{(r^*)^3} \right].
	\]
	
	\textbf{第二步：限制在边界上 ($|x|=a$)}
	当 $x \in \partial B_a$ 时，利用关系式 $r^* = \frac{a}{\rho} r$，我们有：
	\[
	(r^*)^3 = \left( \frac{a}{\rho} \right)^3 r^3 \implies \frac{1}{(r^*)^3} = \frac{\rho^3}{a^3 r^3}.
	\]
	代入梯度的第二项系数：
	\[
	\frac{a}{\rho} \frac{1}{(r^*)^3} = \frac{a}{\rho} \cdot \frac{\rho^3}{a^3 r^3} = \frac{\rho^2}{a^2 r^3}.
	\]
	于是边界上的梯度表达式化简为：
	\[
	\nabla_x G \Big|_{|x|=a} = \frac{1}{4\pi r^3} \left[ (\xi-x) + \frac{\rho^2}{a^2} (x-\xi^*) \right].
	\]
	再将镜像点定义 $\xi^* = \frac{a^2}{\rho^2} \xi$ 代入：
	\[
	\frac{\rho^2}{a^2} (x-\xi^*) = \frac{\rho^2}{a^2} x - \frac{\rho^2}{a^2} \left( \frac{a^2}{\rho^2} \xi \right) = \frac{\rho^2}{a^2} x - \xi.
	\]
	代回梯度公式：
	\[
	\nabla_x G \Big|_{|x|=a} = \frac{1}{4\pi r^3} \left[ \xi - x + \frac{\rho^2}{a^2} x - \xi \right] 
	= \frac{1}{4\pi r^3} \left( \frac{\rho^2}{a^2} - 1 \right) x.
	\]
	
	\textbf{第三步：计算法向导数}
	\[
	\frac{\partial G}{\partial n} = \nabla_x G \cdot \frac{x}{a} 
	= \frac{1}{4\pi r^3} \left( \frac{\rho^2 - a^2}{a^2} \right) x \cdot \frac{x}{a}.
	\]
	由于在球面上 $x \cdot x = |x|^2 = a^2$，故：
	\begin{equation}
		\frac{\partial G}{\partial n} \Bigg|_{|x|=a} 
		= \frac{1}{4\pi r^3} \frac{\rho^2 - a^2}{a^2} \cdot \frac{a^2}{a} 
		= \frac{\rho^2 - a^2}{4\pi a r^3} 
		= - \frac{a^2 - |\xi|^2}{4\pi a |x-\xi|^3}.
	\end{equation}
	
	\paragraph{3. 球上的 Poisson 公式}
	对于 Dirichlet 问题：
	\[
	\begin{cases}
		-\Delta u = 0, & x \in B_a \\
		u = \varphi, & x \in \partial B_a
	\end{cases}
	\]
	解由公式 $u(\xi) = - \int_{\partial B_a} \varphi(x) \frac{\partial G}{\partial n} \mathrm{d}S$ 给出。
	将上述导数结果代入（负负得正），即得\textbf{三维 Poisson 公式}：
	\begin{equation}
		u(\xi) = \frac{a^2 - |\xi|^2}{4\pi a} \int_{\partial B_a} \frac{\varphi(x)}{|x-\xi|^3} \mathrm{d}S_x.
	\end{equation}
	其中 $\xi$ 为球内任意一点。
	
	\begin{note}
		该公式显示，球内任意点的调和函数值是其边界值的加权平均。权函数（Poisson 核）$K(x, \xi) = \frac{a^2 - |\xi|^2}{4\pi a |x-\xi|^3}$ 满足正性 $\frac{a^2 - |\xi|^2}{4\pi a |x-\xi|^3} > 0$（因为 $|\xi| < a$），且在边界上积分归一。
	\end{note}
	
	\subsubsection{高维推广与统一形式*}（可选择用于记忆，可不看）
	
	观察上述二维圆域和三维球域的结果，我们可以发现它们具有高度的内在一致性。这种一致性源于 Laplace 方程基本解的径向对称性以及球几何的特殊性质。
	
	设 $\Omega = B_a(0) \subset \mathbb{R}^n$ 是 $n$ 维空间中的球体 ($n \ge 2$)。记 $\omega_n$ 为 $n$ 维单位球面的表面积（例如 $\omega_2 = 2\pi, \omega_3 = 4\pi$）。
	
	\paragraph{1. 几何结构的统一性}
	无论维数 $n$ 如何，镜像点（反演点）的定义始终为：
	\begin{equation}
		\xi^* = \frac{a^2}{|\xi|^2} \xi, \quad \text{满足 } |\xi| \cdot |\xi^*| = a^2.
	\end{equation}
	且在球面上（$|x|=a$），距离比例关系始终成立：
	\begin{equation}
		\frac{|x-\xi|}{|x-\xi^*|} = \frac{|\xi|}{a}.
	\end{equation}
	
	\paragraph{2. Green 函数的统一表达}
	记 $\Gamma_n(r)$ 为 $n$ 维 Laplace 方程的基本解（仅依赖于距离 $r$）。
	\begin{itemize}
		\item 当 $n=2$ 时，$\Gamma_2(r) = \frac{1}{2\pi} \ln \frac{1}{r}$；
		\item 当 $n \ge 3$ 时，$\Gamma_n(r) = \frac{1}{(n-2)\omega_n} r^{2-n}$.
	\end{itemize}
	利用这一记号，球域上的 Green 函数可以统一写为：
	\begin{equation}
		G(x; \xi) = \Gamma_n(|x-\xi|) - \Gamma_n\left( \frac{|\xi|}{a} |x-\xi^*| \right).
	\end{equation}
	\begin{note}
		上式清晰地表明：Green 函数 $G$ = (源点 $\xi$ 产生的基本解) - (修正后的像点 $\xi^*$ 产生的基本解)。
		修正系数 $\frac{|\xi|}{a}$ 正好使得当 $|x-\xi| = \frac{|\xi|}{a}|x-\xi^*|$（即在边界上）时，两项完全抵消。
	\end{note}
	
	\paragraph{3. Poisson 公式的统一形式}
	对于 Dirichlet 问题，解 $u(\xi)$ 的积分表示为：
	\begin{equation}
		u(\xi) = \int_{\partial B_a} K_n(x, \xi) \varphi(x) \mathrm{d}S_x.
	\end{equation}
	其中 $K_n(x, \xi)$ 称为 \textbf{Poisson 核}。在 $n$ 维情形下，其统一表达式为：
	\begin{equation}
		K_n(x, \xi) = \frac{1}{\omega_n a} \cdot \frac{a^2 - |\xi|^2}{|x-\xi|^n}.
	\end{equation}
	
	\textbf{验证：}
	\begin{itemize}
		\item \textbf{当 n=2 时}：$\omega_2 = 2\pi$，分母为 $|x-\xi|^2$。利用余弦定理展开 $|x-\xi|^2 = a^2 + \rho^2 - 2a\rho \cos\theta$，即得到之前的圆 Poisson 公式。
		\item \textbf{当 n=3 时}：$\omega_3 = 4\pi$，分母为 $|x-\xi|^3$。这正是之前推导出的球 Poisson 公式。
	\end{itemize}
	
	\subsection{球上的泊松积分公式的性质}
	
	\begin{proposition}
		设 $\varphi \in C(\partial B_a(0))$。对于任意 $x \in B_a(0)$，由泊松公式给出的函数 $u$：
		\begin{equation}
			u(x) = \frac{a^2 - |x|^2}{a \omega_n} \int_{\partial B_a(0)} \frac{\varphi(y)}{|x-y|^n} \, d\sigma_y
		\end{equation}
		满足以下性质：
		\begin{enumerate}
			\item $u \in C^\infty(B_a(0))$ 且有界；
			\item $\Delta u = 0$，对于任意 $x \in B_a(0)$；
			\item 对于任意 $x_0 \in \partial B_a(0)$，当 $x \to x_0$ ($x \in B_a(0)$) 时，$u(x) \to \varphi(x_0)$。
		\end{enumerate}
	\end{proposition}
	
	\begin{proof}
		记泊松核为 $K(x, y) = \frac{a^2 - |x|^2}{a \omega_n |x-y|^n}$。
		
		\textbf{1. 证明 $u$ 的光滑性与有界性}
		
		由于 $x \in B_a(0)$ 且 $y \in \partial B_a(0)$，这也是意味着 $|x| < a = |y|$，因此 $|x-y| \neq 0$。
		函数 $K(x, y)$ 关于 $x$ 在 $B_a(0)$ 内是无穷次可微的 (即 $C^\infty$)。根据含参变量积分的求导法则，我们可以将微分算子移入积分号内，从而得到 $u \in C^\infty(B_a(0))$。
		
		关于有界性，已知 $\varphi$ 在紧集 $\partial B_a(0)$ 上连续，故 $\varphi$ 有界。结合泊松核的性质 $\int_{\partial B_a(0)} K(x, y) \, d\sigma_y = 1$，可得：
		\[
		|u(x)| \le \sup_{y \in \partial B_a(0)} |\varphi(y)| \cdot \int_{\partial B_a(0)} K(x, y) \, d\sigma_y = \sup |\varphi| < \infty
		\]
		故 $u$ 有界。
		
		\vspace{1em}
		
		\textbf{2. 证明 $\Delta u = 0$}
		
		只需证明对于固定的 $y \in \partial B_a(0)$，泊松核 $K(x, y)$ 关于 $x$ 是调和的。即验证：
		\[
		\Delta_x \left( \frac{a^2 - |x|^2}{|x-y|^n} \right) = 0
		\]
		计算梯度：
		\begin{align*}
			\nabla_x \left( \frac{a^2 - |x|^2}{|x-y|^n} \right) &= \nabla_x (a^2 - |x|^2) \cdot |x-y|^{-n} + (a^2 - |x|^2) \cdot \nabla_x (|x-y|^{-n}) \\
			&= -2x |x-y|^{-n} + (a^2 - |x|^2) (-n) |x-y|^{-n-2} (x-y)
		\end{align*}
		进一步计算散度 (Laplacian)：
		\begin{align*}
			\Delta_x (\dots) &= \nabla \cdot \left[ -2x |x-y|^{-n} - n(a^2 - |x|^2)|x-y|^{-n-2}(x-y) \right] \\
			&= -2 \left( n |x-y|^{-n} + x \cdot (-n)|x-y|^{-n-2}(x-y) \right) \\
			&\quad - n \left[ \nabla(a^2-|x|^2) \cdot (x-y)|x-y|^{-n-2} + (a^2-|x|^2) \nabla \cdot ((x-y)|x-y|^{-n-2}) \right]
		\end{align*}
		经过代数化简（利用 $|y|=a$ 以及展开 $|x-y|^2 = |x|^2 - 2x \cdot y + |y|^2$），各项恰好抵消，得 $\Delta_x K(x, y) = 0$。因此：
		\[
		\Delta u(x) = \int_{\partial B_a(0)} \varphi(y) \Delta_x K(x, y) \, d\sigma_y = 0
		\]
		
		\vspace{1em}
		
		\textbf{3. 证明边界极限行为}
		
		我们需要证明 $\lim_{x \to x_0} u(x) = \varphi(x_0)$。
		利用性质 $\int_{\partial B_a(0)} K(x, y) \, d\sigma_y = 1$，我们有：
		\[
		u(x) - \varphi(x_0) = \int_{\partial B_a(0)} K(x, y) (\varphi(y) - \varphi(x_0)) \, d\sigma_y
		\]
		对任意 $\epsilon > 0$，由于 $\varphi$ 在 $x_0$ 处连续，存在 $\delta > 0$，使得当 $y \in \partial B_a(0)$ 且 $|y - x_0| < \delta$ 时，有 $|\varphi(y) - \varphi(x_0)| < \frac{\epsilon}{2}$。
		将积分区域 $\partial B_a(0)$ 划分为两部分：$D_1 = \{y : |y - x_0| < \delta\}$ 和 $D_2 = \{y : |y - x_0| \ge \delta\}$。
		\[
		|u(x) - \varphi(x_0)| \le \underbrace{\int_{D_1} K(x, y) |\varphi(y) - \varphi(x_0)| \, d\sigma_y}_{I_1} + \underbrace{\int_{D_2} K(x, y) |\varphi(y) - \varphi(x_0)| \, d\sigma_y}_{I_2}
		\]
		
		对于 $I_1$：
		\[
		I_1 < \frac{\epsilon}{2} \int_{D_1} K(x, y) \, d\sigma_y \le \frac{\epsilon}{2} \int_{\partial B_a(0)} K(x, y) \, d\sigma_y = \frac{\epsilon}{2}
		\]
		
		对于 $I_2$：
		考虑 $x$ 足够接近 $x_0$，即 $|x - x_0| < \frac{\delta}{2}$。
		对于 $y \in D_2$，有 $|y - x_0| \ge \delta$。根据三角不等式：
		\[
		|y - x| \ge |y - x_0| - |x - x_0| > \delta - \frac{\delta}{2} = \frac{\delta}{2}
		\]
		或者更严格地，利用 $|y-x| \ge \frac{1}{2}|y-x_0|$ (当 $|x-x_0|$ 足够小时)。重点在于分母 $|y-x|$ 有下界。
		此时观察 $I_2$ 中的核函数项 $K(x, y)$：
		\[
		K(x, y) = \frac{a^2 - |x|^2}{a \omega_n |x-y|^n}
		\]
		当 $x \to x_0$ 时，分子 $a^2 - |x|^2 \to a^2 - |x_0|^2 = 0$，而分母 $|x-y|^n$ 在 $D_2$ 上有正下界 $(\delta/2)^n$。
		且 $|\varphi(y) - \varphi(x_0)|$ 有界（设上界为 $2M$）。
		因此，存在 $\delta'$ 使得当 $|x - x_0| < \delta'$ 时，
		\[
		I_2 \le 2M \cdot \frac{a^2 - |x|^2}{a \omega_n (\delta/2)^n} \cdot |\partial B_a(0)| < \frac{\epsilon}{2}
		\]
		综上，当 $|x - x_0|$ 足够小时， $|u(x) - \varphi(x_0)| < \epsilon$。证毕。
	\end{proof}
	
	\subsection{极值原理}
	
	本节考虑比 Poisson 方程更一般的二阶椭圆型方程：
	\begin{equation}
		Lu = -\Delta u + c(x)u = f(x). \label{eq:Lu}
	\end{equation}
	设 $\Omega$ 是 $n$ 维欧氏空间 $\mathbb{R}^n$ 的有界开区域。我们在本节中始终假定系数满足非负性条件：
	\begin{equation}
		c(x) \ge 0, \quad \forall x \in \Omega.
	\end{equation}
	这一条件对于最大模估计是至关重要的。
	
	\subsubsection{弱极值原理}
	
	首先引入一个关于严格不等式的基本引理。
	
	\begin{lemma} \label{lemma:strict_max}
		设 $c(x) \ge 0$。若函数 $u \in C^2(\Omega) \cap C(\bar{\Omega})$ 在 $\Omega$ 内满足
		\begin{equation}
			Lu = f < 0,
		\end{equation}
		则 $u$ 不能在 $\Omega$ 内达到它在 $\bar{\Omega}$ 上的\textbf{非负最大值}。
	\end{lemma}
	
	\begin{proof}
		用反证法。假设 $u(x)$ 在 $\Omega$ 内某点 $x^0$ 达到非负最大值，即
		\[
		u(x^0) = \max_{x \in \bar{\Omega}} u(x) = M \ge 0.
		\]
		由多元微积分中极值的必要条件可知，在极大值点 $x^0$ 处：
		\[
		\frac{\partial u}{\partial x_i}(x^0) = 0, \quad \frac{\partial^2 u}{\partial x_i^2}(x^0) \le 0 \quad (i=1, 2, \dots, n).
		\]
		因此 Laplacian 算子满足 $\Delta u(x^0) = \sum \frac{\partial^2 u}{\partial x_i^2}(x^0) \le 0$，即 $-\Delta u(x^0) \ge 0$。
		结合 $c(x) \ge 0$ 和 $u(x^0) \ge 0$，我们有：
		\[
		Lu \big|_{x=x^0} = -\Delta u(x^0) + c(x^0)u(x^0) \ge 0.
		\]
		这与已知条件 $Lu < 0$ 矛盾。引理得证。
	\end{proof}
	
	基于上述引理，我们可以通过扰动方法证明弱极值原理。
	
	\begin{theorem}[弱极值原理] \label{thm:weak_max}
		设 $c(x) \ge 0$ 且在 $\Omega$ 上有界。若 $u \in C^2(\Omega) \cap C(\bar{\Omega})$ 且在 $\Omega$ 内满足 $Lu = f \le 0$，则 $u$ 在 $\bar{\Omega}$ 上的\textbf{非负最大值}必在边界 $\partial \Omega$ 上达到。即：
		\begin{equation}
			\sup_{\Omega} u(x) \le \sup_{x \in \partial \Omega} u^+(x),
		\end{equation}
		其中 $u^+(x) = \max\{ u(x), 0 \}$。
	\end{theorem}
	
	\begin{proof}
		对于任意 $\varepsilon > 0$，构造辅助函数：
		\[
		w(x) = u(x) + \varepsilon e^{\alpha x_1},
		\]
		其中 $\alpha$ 为待定常数。计算 $Lw$：
		\[
		Lw = Lu + \varepsilon L(e^{\alpha x_1}) = Lu + \varepsilon \left( -\frac{\partial^2}{\partial x_1^2} e^{\alpha x_1} + c(x)e^{\alpha x_1} \right).
		\]
		即
		\[
		Lw = Lu + \varepsilon e^{\alpha x_1} (-\alpha^2 + c(x)).
		\]
		由于 $c(x)$ 在 $\Omega$ 内有界，可选取充分大的 $\alpha$，使得对于所有 $x \in \Omega$ 都有 $-\alpha^2 + c(x) < 0$。
		结合 $Lu \le 0$ 及 $\varepsilon > 0$，对于如此选择的 $\alpha$，我们有：
		\[
		Lw < 0.
		\]
		根据引理 \ref{lemma:strict_max}，$w$ 的非负最大值只能在边界 $\partial \Omega$ 上达到。因此：
		\[
		\sup_{\Omega} w(x) \le \sup_{x \in \partial \Omega} w^+(x).
		\]
		于是
		\[
		\sup_{\Omega} u(x) \le \sup_{\Omega} w(x) \le \sup_{x \in \partial \Omega} w^+(x) \le \sup_{x \in \partial \Omega} u^+(x) + \varepsilon \sup_{x \in \Omega} e^{\alpha x_1}.
		\]
		令 $\varepsilon \to 0$，即得 $\sup_{\Omega} u(x) \le \sup_{x \in \partial \Omega} u^+(x)$。
	\end{proof}
	
	\paragraph{注记}
	\begin{enumerate}
		\item 如果 $c(x) \equiv 0$（例如 Laplace 方程），则引理与定理中的“非负最大值”可以用“最大值”代替。
		\item 极值原理具有明显的物理意义：对于一个稳定温度场，如果内部有热汇 ($Lu \le 0$)，那么温度的最大值必在边界达到。
		\item 类似地，如果 $Lu \ge 0$，则 $u$ 在 $\bar{\Omega}$ 上的\textbf{非正最小值}必在 $\partial \Omega$ 上达到。特别地，若 $Lu=0$（调和函数），则最大值与最小值都在边界达到。
	\end{enumerate}
	
	\subsubsection{强极值原理}
	
	定理 2.2 只断言了最大值\textbf{必在}边界达到，但未排除在内部也达到的可能性。强极值原理指出：除非 $u$ 是常数，否则非负最大值\textbf{不能}在内部达到。为证明这一点，首先引入 Hopf 边界点引理。
	
	\begin{lemma}[引理 2.3 : 边界点引理]
		设 $S$ 是 $\mathbb{R}^n$ 中的一个球，在 $S$ 上 $c(x) \ge 0$ 且有界。如果 $u \in C^1(\bar{S}) \cap C^2(S)$ 且满足条件：
		\begin{enumerate}
			\item $Lu \le 0$；
			\item 存在 $x^0 \in \partial S$，使得 $u(x^0) \ge 0$，且当 $x \in S$ 时 $u(x) < u(x^0)$；
		\end{enumerate}
		则有
		\begin{equation}
			\left. \frac{\partial u}{\partial \nu} \right|_{x=x^0} > 0,
		\end{equation}
		其中 $\nu$ 是 $\partial S$ 在点 $x^0$ 处的单位外法向量，且方向导数定义中的方向与 $\nu$ 的夹角小于 $\frac{\pi}{2}$。
	\end{lemma}
	
	利用边界点引理，可以得到如下强极值原理：
	
	\begin{theorem}[定理 2.4 : 强极值原理]
		设 $\Omega$ 是有界连通开区域，在 $\Omega$ 上 $c(x) \ge 0$ 且有界。设 $u \in C(\bar{\Omega}) \cap C^2(\Omega)$ 并满足 $Lu \le 0$。
		如果 $u$ 在 $\Omega$ 内部达到\textbf{非负最大值}，则 $u$ 在 $\Omega$ 上恒为常数。
	\end{theorem}
	
	\paragraph{注记 : 弱极值原理与强极值原理的比较}
	
	在复习时，需明确区分“弱”与“强”的本质差异。强极值原理之所以被称为“强”，是因为它排除了非平凡解在内部达到极值的可能性。
	
	\begin{itemize}
		\item \textbf{条件的差异（区域连通性）}：
		\begin{itemize}
			\item \textbf{弱极值原理}：对区域 $\Omega$ 的拓扑结构无特殊要求（可以是非连通的）。它仅利用了极值点处的局部微分性质（$\Delta u \le 0$）。
			\item \textbf{强极值原理}：必须要求区域 $\Omega$ 是\textbf{连通的}。这是为了利用链式论证（Chaining Argument），将“某一点达到最大值且为常数”的性质延拓到整个区域。
		\end{itemize}
		
		\item \textbf{结论的差异（不等式 vs 恒等式）}：
		\begin{itemize}
			\item \textbf{弱极值原理}：断言 $\displaystyle \sup_{\Omega} u \le \sup_{\partial \Omega} u^+$。它\textbf{不排除} $u$ 在内部达到最大值的可能性（例如常数解）。
			\item \textbf{强极值原理}：断言如果 $u$ 在内部达到非负最大值，则 $u \equiv C$（常数）。换言之，对于非常数解，内部值\textbf{严格小于}边界最大值。
		\end{itemize}
	
	\end{itemize}

	\subsection{边值问题解的最大模估计}
	
	本节主要应用极值原理来建立二阶椭圆型方程边值问题解的先验估计——最大模估计。这类估计在证明解的唯一性及稳定性方面起着核心作用。
	
	\subsubsection{Dirichlet 问题的最大模估计}
	
	考虑 $n$ 维位势方程的 Dirichlet 问题：
	\begin{equation}
		\begin{cases}
			-\Delta u = f(x), & x \in \Omega, \\
			u \big|_{\partial \Omega} = \varphi(x).
		\end{cases} \tag{2.5}
	\end{equation}
	
	\begin{theorem}
		设 $u \in C^2(\Omega) \cap C(\bar{\Omega})$ 是 Dirichlet 问题 (2.5) 的解，则
		\begin{equation}
			\max_{\bar{\Omega}} |u(x)| \le \Phi + C F,
		\end{equation}
		其中 $\Phi = \max_{\partial \Omega} |\varphi(x)|$，$F = \sup_{\Omega} |f(x)|$，$C$ 是仅依赖于维数 $n$ 与 $\Omega$ 的直径的常数。
	\end{theorem}
	
	\begin{proof}
		不妨设区域 $\Omega$ 包含坐标原点 $x=0$。
		令 $d = \sup_{x,y \in \Omega} |x-y|$ 是 $\Omega$ 的直径。
		
		构造辅助函数 $w(x) = z(x) \pm u(x)$，其中 $z(x)$ 选取为：
		\[
		z(x) = \frac{F}{2n}(d^2 - |x|^2) + \Phi.
		\]
		容易验证：
		\[
		-\Delta z = -\frac{F}{2n} \sum_{i=1}^n \frac{\partial^2}{\partial x_i^2}(d^2 - |x|^2) = -\frac{F}{2n} \sum_{i=1}^n (-2) = F.
		\]
		由此可计算 $w$ 的拉普拉斯算子：
		\[
		-\Delta w = -\Delta z \pm (-\Delta u) = F \pm f(x) \ge F - |f(x)| \ge 0.
		\]
		在边界 $\partial \Omega$ 上，由于 $|x| \le d$ 且 $F \ge 0$，有 $z(x) \ge \Phi$。因此：
		\[
		w \big|_{\partial \Omega} = z \big|_{\partial \Omega} \pm \varphi \ge \Phi \pm \varphi \ge \Phi - |\varphi| \ge 0.
		\]
		由极值原理（对于 $-\Delta w \ge 0$），$w(x)$ 的非负最小值必在边界达到，故在 $\Omega$ 上 $w(x) \ge 0$。
		这意味着 $z(x) \pm u(x) \ge 0$，即 $|u(x)| \le z(x)$。
		
		注意到 $|x| \ge 0$，我们有：
		\[
		|u(x)| \le \frac{F}{2n} d^2 + \Phi = \Phi + \frac{d^2}{2n} F.
		\]
		取 $C = \frac{d^2}{2n}$，定理证毕。
	\end{proof}
	
	\subsubsection{第三边值问题的最大模估计}
	
	在 $\Omega$ 上考虑第三边值问题：
	\begin{equation}
		\begin{cases}
			Lu = -\Delta u + c(x)u = f(x), & x \in \Omega, \\
			\left( \frac{\partial u}{\partial n} + \alpha(x)u \right) \bigg|_{\partial \Omega} = \varphi(x).
		\end{cases} \tag{2.6}
	\end{equation}
	其中 $n$ 表示 $\partial \Omega$ 的单位外法向。如果 $\alpha(x) \equiv 0$，则问题 (2.6) 称为 Neumann 问题或第二边值问题。
	
	\begin{theorem}
		设 $c(x) \ge 0$，$\alpha(x) \ge \alpha_0 > 0$，$u \in C^2(\Omega) \cap C^1(\bar{\Omega})$ 是问题 (2.6) 的解，则
		\begin{equation}
			\max_{\bar{\Omega}} |u(x)| \le C(\Phi + F),
		\end{equation}
		其中 $C$ 只依赖于 $n, \alpha_0, d$，而 $\Phi = \max_{\partial \Omega}|\varphi(x)|$，$F = \sup_{\Omega}|f(x)|$。
	\end{theorem}
	
	\begin{proof}
		不妨设区域包含坐标原点 $x=0$。考虑辅助函数
		\[
		w(x) = z(x) + \frac{\Phi}{\alpha_0} \pm u(x),
		\]
		其中选取
		\[
		z(x) = \frac{F}{2n} \left( \frac{1+d^2}{\alpha_0} + d^2 - |x|^2 \right).
		\]
		记 $\partial \Omega$ 的单位外法向 $n = (\beta_1(x), \beta_2(x), \dots, \beta_n(x))$，则 $\sum_{i=1}^n \beta_i^2(x) = 1$ 在 $x \in \partial \Omega$ 上成立。
		
		首先验证方程：
		\[
		-\Delta z = F \implies -\Delta z + c(x)z \ge F \quad (\text{因 } c \ge 0, z \ge 0).
		\]
		于是
		\begin{equation}
			-\Delta w + c(x)w = (-\Delta z + cz) + c\frac{\Phi}{\alpha_0} \pm (-\Delta u + cu) \ge F \pm f \ge 0. \tag{2.7}
		\end{equation}
		
		其次验证边界条件。计算 $z$ 的边界算子：
		\[
		\frac{\partial z}{\partial n} = \nabla z \cdot n = \frac{F}{2n} (-2x) \cdot n = -\frac{F}{n} \sum_{i=1}^n x_i \beta_i(x).
		\]
		于是
		\[
		\begin{aligned}
			\left[ \frac{\partial z}{\partial n} + \alpha(x)z \right]_{\partial \Omega} 
			&= \frac{F}{2n} \left[ -2\sum_{i=1}^n x_i \beta_i + \alpha(x) \left( \frac{1+d^2}{\alpha_0} + d^2 - |x|^2 \right) \right]_{\partial \Omega} \\
			&\ge \frac{F}{2n} \left[ -|x|^2 - \sum \beta_i^2(x) + \alpha_0 \frac{1+d^2}{\alpha_0} \right]_{\partial \Omega} \quad (\text{利用不等式 } 2|A||B| \le A^2+B^2) \\
			&\ge \frac{F}{2n} \left[ -|x|^2 - 1 + 1 + d^2 \right] = \frac{F}{2n}(d^2 - |x|^2) \ge 0.
		\end{aligned}
		\]
		因此，对于 $w$ 的边界条件：
		\begin{equation}
			\begin{aligned}
				\left[ \frac{\partial w}{\partial n} + \alpha(x)w \right]_{\partial \Omega} 
				&= \left[ \frac{\partial z}{\partial n} + \alpha z \right] + \alpha(x)\frac{\Phi}{\alpha_0} \pm \varphi(x) \\
				&\ge 0 + \frac{\alpha(x)}{\alpha_0}\Phi - |\varphi(x)| \ge \Phi - \Phi = 0.
			\end{aligned} \tag{2.8}
		\end{equation}
		
		\textbf{结论推导：}
		由弱极值原理（参见附注3关于非正最小值），$w(x)$ 若有负最小值，必在边界 $\partial \Omega$ 上达到。
		假设 $w$ 在某点 $x^0 \in \partial \Omega$ 达到\textbf{负最小值}（即 $w(x^0) < 0$），则在该点处必有 $\frac{\partial w}{\partial n} \le 0$（外法向导数）。
		结合 $w(x^0) < 0$ 和 $\alpha(x) > 0$，有：
		\[
		\left[ \frac{\partial w}{\partial n} + \alpha(x)w \right]_{x=x^0} \le \alpha(x^0) w(x^0) < 0.
		\]
		这与已证的边界不等式 (2.8) 矛盾！
		这就说明 $w(x)$ 在 $\bar{\Omega}$ 上不可能有负最小值，于是 $w(x) \ge 0$。
		即
		\[
		|u(x)| \le z(x) + \frac{\Phi}{\alpha_0} \le \frac{F}{2n} \left( \frac{1+d^2}{\alpha_0} + d^2 \right) + \frac{\Phi}{\alpha_0}.
		\]
		整理得
		\[
		\max_{\bar{\Omega}} |u(x)| \le C(F + \Phi),
		\]
		其中
		\[
		C = \max \left\{ \frac{1}{\alpha_0}, \frac{1}{2n} \left( \frac{1+d^2}{\alpha_0} + d^2 \right) \right\}.
		\]
	\end{proof}
	
	
	\subsection{调和函数的性质}
	
	在研究 Laplace 方程之前，我们首先给出调和函数的定义。调和函数不仅是 Laplace 方程的解，更具有许多优美的解析性质。
	
	\begin{definition}[调和函数]
		设 $\Omega$ 是 $\mathbb{R}^n$ 中的开区域。如果函数 $u \in C^2(\Omega)$ 满足 Laplace 方程：
		\begin{equation}
			\Delta u = 0, \quad \forall x \in \Omega
		\end{equation}
		则称 $u$ 是 $\Omega$ 内的\textbf{调和函数}。
	\end{definition}
	
	以下是调和函数的一些核心性质：
	
	\subsubsection{解析性}
	
	调和函数最惊人的性质之一是它的正则性。尽管定义只要求 $C^2$，但实际上它具有最高阶的光滑性。
	
	\begin{theorem}[定理 : 解析性]
		设 $u$ 在 $\Omega$ 内调和，则 $u$ 在 $\Omega$ 内解析。
	\end{theorem}
	
	\begin{note}
		这意味着调和函数是 $C^\omega(\Omega)$ 类函数（即在每一点邻域内可展成幂级数），当然也是 $C^\infty(\Omega)$ 的。这表明 Laplace 方程具有极强的\textbf{平滑效应}：即使边界值只是连续的，内部解也是无穷次可微的。
	\end{note}
	
	\subsubsection{平均值定理}
	
	平均值性质是调和函数的几何特征，它也是极值原理的直观来源。
	
	\begin{theorem}[定理 : 平均值定理]
		设 $u(x,y)$ 在 $\Omega$ 内调和。对于任意点 $P_0(x_0, y_0) \in \Omega$，如果以 $P_0$ 为心、半径为 $R$ 的闭圆 $\bar{B}_R(P_0) \subset \Omega$，则有：
		\begin{equation}
			u(x_0, y_0) = \frac{1}{2\pi R} \int_{\partial B_R(P_0)} u \, \mathrm{d}l.
		\end{equation}
	\end{theorem}
	
	\begin{note}
		该定理表明，调和函数在圆心的值等于其在圆周上的平均值。
		在三维情形下，类似地有球平均公式：
		\[
		u(x_0) = \frac{1}{4\pi R^2} \int_{\partial B_R(x_0)} u(x) \, \mathrm{d}S.
		\]
		事实上，不仅是在球面上，调和函数在球体\textbf{内部}的平均值也等于球心的值。
	\end{note}
	
	\subsubsection{Liouville 定理}
	
	Liouville 定理刻画了全空间上调和函数的刚性。它是复变函数中 Liouville 定理在实调和函数中的推广。
	
	\begin{theorem}[定理 : Liouville 定理]
		在全平面上有下界（或有上界）的调和函数必为常数。
	\end{theorem}
	
	\begin{proof}[思路简述]
		通常利用梯度估计（如 Bernstein 估计）或平均值定理证明。若 $u$ 有界，可以推导出 $|\nabla u| \equiv 0$，从而 $u$ 为常数。
		值得注意的是，对于调和函数，只需要\textbf{单侧有界}（例如非负 $u \ge 0$）即可推出常数结论，这比复变函数中的要求（模有界）更强。
	\end{proof}
	
	\section{重要习题}
	
	\subsection*{习题 1：无界区域上的极值原理}
	
	\textbf{题目：}
	设 $\Omega = \{ x \in \mathbb{R}^n \mid |x| > 1 \}$ 为单位球外部的无界区域。
	设 $u \in C^2(\Omega) \cap C(\bar{\Omega})$ 是下列边值问题的解：
	\begin{equation}
		\begin{cases}
			\Delta u = 0, & x \in \Omega \\
			\lim_{|x| \to \infty} u(x) = 0. &
		\end{cases}
	\end{equation}
	证明：$u$ 在 $\bar{\Omega}$ 上的最大模在边界 $\partial \Omega$（即 $|x|=1$）上达到，即
	\begin{equation}
		\max_{\bar{\Omega}} |u(x)| = \max_{\partial \Omega} |u(x)|.
	\end{equation}
	
	\begin{proof}[\textbf{证明}]
		由于区域 $\Omega$ 无界，不能直接应用有界区域的极值原理。我们采用截断区域法。
		对于任意给定的 $\varepsilon > 0$，由无穷远处的衰减条件 $\lim_{|x| \to \infty} u(x) = 0$ 可知，存在充分大的 $R > 1$，使得当 $|x| \ge R$ 时，
		\[
		|u(x)| < \varepsilon.
		\]
		现在考虑有界圆环区域 $D_R = \{ x \in \mathbb{R}^n \mid 1 < |x| < R \}$。
		$u$ 在 $D_R$ 上调和，由有界区域的最大模估计（定理 2.5 或推论），$u$ 在 $\bar{D}_R$ 上的最大模必在边界 $\partial D_R$ 上达到。
		$\partial D_R$ 由两部分组成：内边界 $\partial \Omega$ ($|x|=1$) 和外边界 $S_R$ ($|x|=R$)。
		因此：
		\[
		\max_{\bar{D}_R} |u(x)| = \max \{ \max_{|x|=1} |u(x)|, \max_{|x|=R} |u(x)| \}.
		\]
		由此可得，对于任意 $x \in D_R$：
		\[
		|u(x)| \le \max \{ \max_{\partial \Omega} |u|, \varepsilon \}.
		\]
		固定 $x$，令 $R \to \infty$（注意不等式对充分大的 $R$ 恒成立），得：
		\[
		|u(x)| \le \max \{ \max_{\partial \Omega} |u|, \varepsilon \}.
		\]
		再令 $\varepsilon \to 0$，即得 $|u(x)| \le \max_{\partial \Omega} |u|$。
		证毕。
	\end{proof}
	
	---
	
	\subsection*{习题 2：波动方程的球面平均 (Darboux 方程)}
	
	\textbf{题目：}
	设 $u(x,t) \in C^2(\mathbb{R}^n \times [0, \infty))$ 是波动方程
	\[
	u_{tt} - c^2 \Delta u = 0
	\]
	的解。定义 $u$ 在以 $x$ 为心、半径为 $r$ 的球面上的平均值（球面平均）为：
	\begin{equation}
		M_u(x, r, t) = \frac{1}{\omega_n} \int_{|\xi|=1} u(x + r\xi, t) \mathrm{d}\sigma_\xi,
	\end{equation}
	其中 $\omega_n$ 是 $n$ 维单位球面的面积。
	证明 $M_u$ 满足 Euler-Poisson-Darboux 方程：
	\begin{equation}
		\frac{\partial^2}{\partial t^2} M_u = c^2 \left( \frac{\partial^2}{\partial r^2} + \frac{n-1}{r} \frac{\partial}{\partial r} \right) M_u.
	\end{equation}
	
	\begin{proof}[\textbf{证明}]
		\textbf{第一步：计算 $M_u$ 关于时间 $t$ 的二阶导数}
		
		由于 $t$ 仅出现在被积函数 $u$ 中，且积分区域（单位球面 $|\xi|=1$）与 $t$ 无关，我们可以直接在积分号下求导：
		\[
		\frac{\partial^2}{\partial t^2} M_u(x, r, t) = \frac{1}{\omega_n} \int_{|\xi|=1} u_{tt}(x + r\xi, t) \, \mathrm{d}\sigma_\xi.
		\]
		利用波动方程 $u_{tt} = c^2 \Delta_y u$（此处记 $y = x+r\xi$ 为空间变量），代入上式得：
		\begin{equation} \label{eq:mut_laplacian}
			\frac{\partial^2}{\partial t^2} M_u(x, r, t) = \frac{c^2}{\omega_n} \int_{|\xi|=1} \Delta_y u(x + r\xi, t) \, \mathrm{d}\sigma_\xi.
		\end{equation}
		
		\textbf{第二步：计算 $M_u$ 关于半径 $r$ 的导数}
		
		对 $M_u$ 关于 $r$ 求一阶导数。注意 $r$ 出现在 $u$ 的自变量中，利用链式法则 $\frac{\partial}{\partial r} u(x+r\xi) = \nabla u(x+r\xi) \cdot \xi$：
		\[
		\frac{\partial}{\partial r} M_u(x, r, t) = \frac{1}{\omega_n} \int_{|\xi|=1} \nabla u(x + r\xi, t) \cdot \xi \, \mathrm{d}\sigma_\xi.
		\]
		注意到 $\xi$ 正是球面 $\partial B(x,r)$ 在点 $y=x+r\xi$ 处的单位外法向量 $\nu$。我们将上述积分变换到以 $x$ 为心、半径为 $r$ 的球体 $B(x,r)$ 上。
		
		首先，将单位球面上的积分转化为半径为 $r$ 的球面 $\partial B(x,r)$ 上的积分（面积元关系为 $\mathrm{d}S_y = r^{n-1}\mathrm{d}\sigma_\xi$）：
		\[
		\int_{|\xi|=1} \nabla u \cdot \xi \, \mathrm{d}\sigma_\xi = \frac{1}{r^{n-1}} \int_{\partial B(x,r)} \frac{\partial u}{\partial \nu} \, \mathrm{d}S_y.
		\]
		接下来，应用散度定理：
		\[
		\int_{\partial B(x,r)} \frac{\partial u}{\partial \nu} \, \mathrm{d}S_y = \int_{B(x,r)} \Delta u(y, t) \, \mathrm{d}y.
		\]
		结合以上两式，我们得到一阶径向导数的表达式：
		\begin{equation} \label{eq:mur_first}
			\frac{\partial}{\partial r} M_u(x, r, t) = \frac{1}{\omega_n r^{n-1}} \int_{B(x,r)} \Delta u(y, t) \, \mathrm{d}y.
		\end{equation}
		
		\textbf{第三步：构建 Euler-Poisson-Darboux 方程右端项}
		
		将方程 \eqref{eq:mur_first} 变形为：
		\[
		r^{n-1} \frac{\partial}{\partial r} M_u = \frac{1}{\omega_n} \int_{B(x,r)} \Delta u(y, t) \, \mathrm{d}y.
		\]
		对上式两边关于 $r$ 求导。注意右边是对变限区域 $B(x,r)$ 的积分求导，这等价于在边界 $\partial B(x,r)$ 上积分（或将体积元写为极坐标形式 $\mathrm{d}y = \rho^{n-1}\mathrm{d}\sigma_\xi \mathrm{d}\rho$ 后对上限求导）：
		\[
		\frac{\partial}{\partial r} \left( r^{n-1} \frac{\partial}{\partial r} M_u \right) = \frac{1}{\omega_n} \int_{\partial B(x,r)} \Delta u(y, t) \, \mathrm{d}S_y.
		\]
		再次利用面积元变换 $\mathrm{d}S_y = r^{n-1}\mathrm{d}\sigma_\xi$ 将积分域变回单位球面：
		\[
		\frac{\partial}{\partial r} \left( r^{n-1} \frac{\partial}{\partial r} M_u \right) = \frac{1}{\omega_n} \int_{|\xi|=1} \Delta u(x+r\xi, t) \, r^{n-1} \mathrm{d}\sigma_\xi.
		\]
		消去两边的 $r^{n-1}$（假设 $r > 0$）：
		\[
		\frac{1}{r^{n-1}} \frac{\partial}{\partial r} \left( r^{n-1} \frac{\partial}{\partial r} M_u \right) = \frac{1}{\omega_n} \int_{|\xi|=1} \Delta u(x+r\xi, t) \, \mathrm{d}\sigma_\xi.
		\]
		
		\textbf{第四步：综合}
		
		观察上式右端，它正是方程 \eqref{eq:mut_laplacian} 右端除以 $c^2$ 的部分。因此我们有：
		\[
		\frac{1}{c^2} \frac{\partial^2}{\partial t^2} M_u = \frac{1}{r^{n-1}} \frac{\partial}{\partial r} \left( r^{n-1} \frac{\partial}{\partial r} M_u \right).
		\]
		展开左边的微分算子：
		\[
		\frac{1}{r^{n-1}} \left( (n-1)r^{n-2} \frac{\partial M_u}{\partial r} + r^{n-1} \frac{\partial^2 M_u}{\partial r^2} \right) = \frac{\partial^2 M_u}{\partial r^2} + \frac{n-1}{r} \frac{\partial M_u}{\partial r}.
		\]
		整理即得 Euler-Poisson-Darboux 方程：
		\[
		\frac{\partial^2}{\partial t^2} M_u = c^2 \left( \frac{\partial^2}{\partial r^2} + \frac{n-1}{r} \frac{\partial}{\partial r} \right) M_u.
		\]
	\end{proof}
	
	
	\subsection*{习题 3：热传导方程初值问题的解}
	
	\textbf{题目：}
	考虑一维热传导方程的 Cauchy 问题：
	\begin{equation}
		\begin{cases}
			u_t - a^2 u_{xx} = 0, & x \in \mathbb{R}, t > 0 \\
			u(x, 0) = \varphi(x). &
		\end{cases}
	\end{equation}
	设 $\varphi(x)$ 连续且有界。利用热核（基本解）$K(x,t) = \frac{1}{\sqrt{4\pi a^2 t}} e^{-\frac{x^2}{4a^2 t}}$，解可表示为：
	\begin{equation}
		u(x,t) = \int_{-\infty}^{+\infty} K(x-\xi, t) \varphi(\xi) \mathrm{d}\xi.
	\end{equation}
	验证上述表达式确实是该初值问题的解。即证明：
	1. $u(x,t)$ 满足方程 $u_t = a^2 u_{xx}$ ($t>0$)；
	2. $\lim_{t \to 0^+} u(x,t) = \varphi(x)$.
	
	\begin{proof}
		\textbf{第一步：验证 $u(x,t)$ 满足热传导方程 ($t>0$)}
		
		当 $t>0$ 时，由于热核 $K(x-\xi, t)$ 具有极好的光滑性（无穷次可微），且 $\varphi(\xi)$ 有界，我们可以直接在积分号下对参数 $x$ 和 $t$ 求导。
		注意到热核 $K$ 本身就是热传导方程的基本解，即满足：
		\[
		\frac{\partial K}{\partial t} - a^2 \frac{\partial^2 K}{\partial x^2} = 0
		\]
		因此：
		\begin{align*}
			\frac{\partial u}{\partial t} - a^2 \frac{\partial^2 u}{\partial x^2} &= \int_{-\infty}^{\infty} \left( \frac{\partial}{\partial t} - a^2 \frac{\partial^2}{\partial x^2} \right) K(x-\xi, t) \varphi(\xi) \mathrm{d}\xi \\
			&= \int_{-\infty}^{\infty} 0 \cdot \varphi(\xi) \mathrm{d}\xi = 0
		\end{align*}
		这说明 $u(x,t)$ 在 $t>0$ 时确实满足齐次热传导方程。
		
		\vspace{0.5em}
		
		\textbf{第二步：验证初始条件 ($t \to 0^+$)}
		
		我们需要证明对于任意 $x_0 \in (-\infty, \infty)$，有 $\lim_{x \to x_0, t \to 0^+} u(x,t) = \varphi(x_0)$。
		
		直接在原积分中取 $t \to 0$ 比较困难，因为 $K$ 会变成狄拉克 $\delta$ 函数。我们采用变量代换法，将 $t$ 从指数部分移到函数自变量中。
		令积分变量变换：
		\[
		\eta = \frac{\xi - x}{2a\sqrt{t}} \quad \Rightarrow \quad \xi = x + 2a\sqrt{t}\eta, \quad \mathrm{d}\xi = 2a\sqrt{t}\mathrm{d}\eta
		\]
		代入 $u(x,t)$ 的表达式，系数消去，积分上下限保持 $-\infty$ 到 $+\infty$：
		\begin{equation}
			u(x,t) = \frac{1}{\sqrt{\pi}} \int_{-\infty}^{\infty} \mathrm{e}^{-\eta^2} \varphi(x + 2a\sqrt{t}\eta) \mathrm{d}\eta
		\end{equation}
		现在我们可以对上式取极限。由于 $\varphi$ 有界，根据实变函数中的控制收敛定理（或利用一致收敛性），可以将极限运算放入积分号内部：
		\begin{align*}
			\lim_{x \to x_0, t \to 0^+} u(x,t) &= \frac{1}{\sqrt{\pi}} \int_{-\infty}^{\infty} \lim_{x \to x_0, t \to 0^+} \left[ \mathrm{e}^{-\eta^2} \varphi(x + 2a\sqrt{t}\eta) \right] \mathrm{d}\eta \\
			&= \frac{1}{\sqrt{\pi}} \int_{-\infty}^{\infty} \mathrm{e}^{-\eta^2} \varphi(x_0) \mathrm{d}\eta
		\end{align*}
		将常数 $\varphi(x_0)$ 提到积分号外，利用高斯积分公式 $\int_{-\infty}^{\infty} \mathrm{e}^{-\eta^2} \mathrm{d}\eta = \sqrt{\pi}$，我们得到：
		\[
		\lim_{x \to x_0, t \to 0^+} u(x,t) = \varphi(x_0) \cdot \frac{1}{\sqrt{\pi}} \cdot \sqrt{\pi} = \varphi(x_0)
		\]
		证毕。
	\end{proof}
	
	
\end{document}